% Independent R6 audit, statements E.110 onward.
\subsection{The actual elliptic cover and its arithmetic}

\begin{result}{E.110}{Proved}{R6 C.3.1: the pointed elliptic cover and the necessary factorization}
Let $Y\to U$ be the pointed degree-$27$ cover defined by the
exponent-three Heisenberg quotient and the section of $t=3$. Its
compactification, with the distinguished point above $3$ as origin, is
$E':y^2=x^3+9$. Put $E:y^2=x^3-243$, $P_0=(6,15)$, and $T=(0,3)$ on $E'$.
The map is
\[
 E'\xrightarrow{\phi}E\xrightarrow{\iota^{-1}}
 C:\ 3u^3-2v^3=1\xrightarrow{3u^3}\mathbb P^1,
\]
where, for $d=3u(u-1)-2v(v-1)$,
\[
 \phi(x,y)=\left(\frac{x^3+36}{x^2},
                    \frac{y(x^3-72)}{x^3}\right),\qquad
 \iota(u,v)=\left(\frac{18(uv-1)}d,
       -\frac{27(3u(u+1)+2v(v+1))}d\right).
\]
There are nine punctures above each of $0,1,\infty$. Every full
geometrically pro-three section above an original A.44 pair factors,
after geometric conjugacy, through the actual geometrically pro-three
arithmetic group of this $Y$.
\end{result}
\begin{proof}
The central quotient is the pointed cubic $C$, with origin $(1,1)$,
by the two Kummer equations $t=3u^3$, $1-t=-2v^3$. Here is a proof that
the displayed $\iota$ is an isomorphism which does not assume the
previous run's inverse formula. Substitution gives $y^2=x^3-243$.
Moreover,
\[
 \operatorname{Res}_v(3u^3-2v^3-1,d)
       =12(u-1)^3(3u^3+2).
\]
At $P=(1,1)$, with local parameter $z=u-1$, the cubic equation gives
\[
 v=1+\tfrac32z-\tfrac34z^2+\tfrac{13}{8}z^3+O(z^4),\quad
 d=\tfrac54z^3+O(z^4),\quad uv-1=\tfrac52z+O(z^2).
\]
Thus $x\circ\iota$ has a pole of order two at $P$.
The other three denominator zeros are $uv=1$, $u^3=-2/3$;
they are simple and both coordinate numerators vanish there.
There is no denominator zero at infinity: $3U^3=2V^3$ and
$3U^2=2V^2$ would force $U=V=0$. Hence the pole divisor of
$x\circ\iota$ is exactly $2P$. The rational map to the smooth
projective curve $E$ extends everywhere. Since $x:E\to\mathbb P^1$
has degree two, the map $\iota$ has degree one. It sends $P$ to
$O$, so it is a pointed isomorphism.

The three inertia subgroups in the Heisenberg group have order three,
and their images in its abelianization are nonzero. Riemann--Hurwitz
gives $2g(\overline Y)-2=-54+3\cdot18=0$; the center meets every
inertia subgroup trivially. Thus the central quotient is an unramified
pointed isogeny $f:\overline Y\to E$ of degree three.
The displayed $\phi$ has kernel $\{O,T,-T\}$. Its dual is
\[
 \psi(x,y)=\left(\frac{x^3-972}{9x^2},
                 \frac{y(x^3+1944)}{27x^3}\right).
\]
The chord-and-tangent formula gives $\phi\psi=[3]$; surjectivity
then also gives $\psi\phi=[3]$. The three-division polynomial of $E$
is $3x(x^3-972)$. Since $972$ is not a rational cube, $E$ has just
one rational subgroup scheme of order three. Dualizing $f$, its dual
has the same kernel as $\psi$. The factorization property of isogenies
therefore identifies $\overline Y$ with $E'$ and $f$ with $\phi$.
The precise elementary isogeny facts used here are
\citet[II.1.5--1.8, p.~49]{AuditMilneEC}. The degree-$27$ map has
ramification index three at every boundary point, proving the boundary
count. All polynomial identities above were recalculated in
\texttt{code/ghost\_audit\_elliptic.py}.

We give the section argument as well, since identifying a curve does
not identify the subgroup through which a section must factor.
Both A.44 coordinates reduce modulo cubes to $(3,-2)$: both new
exponents are $n_j=(j+4)!$, including $n_0=24$. Normalize literal
abelian cocycles to those of the point-three section modulo three.
The difference of two class-two central cochains is then an
$\mathbb F_3(2)=\mathbb F_3$ cocycle. At every new support prime,
the integral tame relation and the ordinary unit equation give
central inertia $n_j u$ or $-n_jv$; it is zero modulo three.
At other primes outside $\{2,3\}$ both inertia coordinates are zero
and the integral tame relation forces central inertia zero. The
central difference is therefore unramified outside $\{2,3\}$.
Kronecker--Weber identifies this character space with the unique
cubic character of conductor nine, generated by
$\eta=(\chi^2-1)/24\bmod3$; powers of two contribute no cubic
character. The reduction has the form $h_3+k\eta$, $k\in\mathbb F_3$.

At the rational prime two the original pair has exactly the same
three-primary local Kummer classes as $(3,-2)$. Its local abelian
alignment can be made by an integral conjugator divisible by three:
$\chi(F_2)-1=1$ is a three-adic unit. The third obstruction at two
changes by $k\eta(F_2)=k/8$ in $\mathbb Z_3/(2^2-1)$, as in B.101.
The point-three section has zero obstruction, so an integral class-three
section must have $k=0$. A full pro-three section therefore projects
to the point-three section in $H_3$ and lies in the subgroup defining
$Y$. A subgroup of three-power index in a pro-three group has its
full pro-three topology: intersect the finitely many conjugates of
an open normal subgroup to obtain an ambient normal subgroup still
of three-power index. This proves factorization through the actual
pro-three arithmetic group of $Y$, not merely a quotient of it.
\end{proof}

\begin{result}{E.111}{Proved}{R6 C.3.2: a complete three-isogeny descent}
With the domain curve in the Selmer notation,
\[
 \operatorname{Sel}^{\phi}(E'/\mathbb Q)=0,\qquad
 \operatorname{Sel}^{\psi}(E/\mathbb Q)
      =\langle6,12\rangle\simeq(\mathbb Z/3)^2.
\]
Both curves have zero three-primary Tate--Shafarevich group, and
\[
 E'(\mathbb Q)^{\wedge3}
     =\mathbb Z_3P_0\oplus(\mathbb Z/3)T.
\]
In particular, the rank is one. This does not assert that $P_0$ is
an integral generator of the entire Mordell--Weil group; it is a
generator of its free three-adic completion.
\end{result}
\begin{proof}
The two kernels are respectively the constant group $\mathbb Z/3$
and $\mu_3$, since the nonzero kernel points of $\psi$ are
$(0,\pm9\sqrt{-3})$. Their discriminants are supported on two and
three. Away from those primes the isogenies extend to finite etale
maps of the smooth proper models, so the local Kummer images are
unramified. Thus every Selmer class is unramified outside $\{2,3\}$.

For the constant kernel, Kronecker--Weber
\citep[I.4.16, pp.~49--50]{RootMilneCFT} leaves a one-dimensional
candidate space: the unique cubic character of conductor nine.
There are no nonzero rational three-torsion points on $E$ over
$\mathbb Q_2$. Indeed its division polynomial permits only $x=0$,
which would require a square root of $-243\equiv5\pmod8$, or
$x^3=972$, impossible because $v_2(972)=2$. An open subgroup of
$E(\mathbb Q_2)$ is a pro-two formal group and is uniquely
three-divisible. Its finite quotient has no three-torsion: otherwise
one corrects a lift by unique division in that open subgroup and
obtains actual three-torsion. Consequently $[3]$ is an automorphism
of $E(\mathbb Q_2)$. Since $\phi\psi=[3]$, $\phi$ is locally
surjective. The nonzero global candidate characters have nonzero
value $\eta(F_2)=1/8\pmod3$, so neither lies in the local Kummer
image. This proves $\operatorname{Sel}^{\phi}=0$.

For the cyclotomic kernel, the candidate space consists of the nine
cube classes $2^a3^b$, $a,b\in\mathbb F_3$. The actual isogeny
Kummer function is $Y-3$ on $E'$. To check its identification, not
merely its divisor, direct substitution gives
\[
 (Y-3)\circ\psi=\left(\frac{y-27}{3x}\right)^3.
\]
Conversely, if $r^3=Y-3$ and $s=X/r$, then $s^3=Y+3$ and
$x=18/(s-r)$, $y=27(s+r)/(s-r)$ recover the fibre. Hence this is
precisely its Kummer torsor. The rational points $-T$ and $P_0$
give the classes $-6=6$ and $12$. Their valuation vectors at
$(2,3)$ are $(1,1)$ and $(2,1)$, independent over $\mathbb F_3$.
All candidate classes are therefore realized by global points;
this proves every local condition, including those at two and three,
and proves $\operatorname{Sel}^{\psi}=\langle6,12\rangle$.

The two exact isogeny Selmer sequences give
$\operatorname{Sha}(E')[\phi]=0$ and
$\operatorname{Sha}(E)[\psi]=0$, as well as
$E(\mathbb Q)=\phi E'(\mathbb Q)$. Composing the isogenies
shows that both groups have zero three-torsion and hence zero
three-primary torsion. Also $\psi E(\mathbb Q)=3E'(\mathbb Q)$,
so the quotient $E'(\mathbb Q)/3$ has order nine, generated by
$P_0$ and $T$. Its rational three-torsion is exactly $\langle T\rangle$,
as follows from $3x(x^3+36)$. The finite-basis theorem
\citep[Chapter IV, p.~105]{AuditMilneEC} now gives rank one.
There is no rational nine-torsion: it would make $T$ divisible by
three, whereas its Kummer class modulo $\psi$ is nonzero.
The class of $P_0$ is nonzero in the free quotient modulo three;
successive three-adic digit lifting yields the stated completion.
The isogeny Selmer sequences used here follow directly from the
isogeny Galois-module exact sequences and their localizations, just
as the multiplication sequence in
\citet[IV.2, (28)--(29)]{AuditMilneEC}; they involve no numerical
Selmer guess or unproved rank computation.
\end{proof}

\begin{result}{E.112}{Proved}{R6 C.3.3 and a stronger localization statement at three}
The natural maps are isomorphisms
\[
 E'(\mathbb Q)^{\wedge3}
 \xrightarrow{\sim}H^1_{\rm cts}(G_{\mathbb Q},T_3E'),
 \qquad
 E'(\mathbb Q)^{\wedge3}
 \xrightarrow{\sim}E'(\mathbb Q_3)^{\wedge3}.
\]
The first assertion is absolute continuous cohomology, without a
fixed ramification set. The second strengthens the previous audit:
the compact global class is detected exactly at the coefficient prime
three. Moreover $\phi:E'(\mathbb Q_3)\to E(\mathbb Q_3)$ is
surjective.
\end{result}
\begin{proof}
First prove the local assertion. The compact three-adic Lie group
$E'(\mathbb Q_3)$ has an open subgroup isomorphic to $\mathbb Z_3$.
Its three-torsion has order three: $x^3=-36$ is impossible over
$\mathbb Q_3$, since its valuation is two. Thus its quotient by
three has order $3\cdot3=9$. The cube classes $6$ and $12$ are
independent even over $\mathbb Q_3$. If $6^a12^b$ is a cube, its
valuation gives $a+b=0\pmod3$, and its unit part $2^{a+2b}$ must
be a cube. Unit cubes modulo nine are $1,-1$, so the class of $2$
has order three and $a+2b=0\pmod3$. Hence $a=b=0$.
The $\psi$-Kummer classes of $-T,P_0$ therefore already generate
a subgroup of order nine in $E'(\mathbb Q_3)/\psi E(\mathbb Q_3)$.
Since $3E'(\mathbb Q_3)\subseteq\psi E(\mathbb Q_3)$, this
quotient has exactly order nine and equals the quotient by three.
It follows that $\psi E(\mathbb Q_3)=3E'(\mathbb Q_3)$.
The kernel of $\psi$ has no nonzero $\mathbb Q_3$-point, because
$-3$ is not a square there; hence
$E(\mathbb Q_3)=\phi E'(\mathbb Q_3)$.

Also $T$ is not divisible by three locally, since its $\psi$-Kummer
class is nonzero. Thus the entire local three-primary torsion is
$\langle T\rangle$. The two rational points generate the local
quotient modulo three, so generate the local pro-three completion
by digit lifting. The map
$\mathbb Z_3P_0\oplus(\mathbb Z/3)T\to E'(\mathbb Q_3)^{\wedge3}$
is surjective between modules of the same rank, and is injective:
the free generator maps to a generator of the free quotient, and
$T$ maps nontrivially. This proves the second isomorphism.

For the first assertion, the finite Kummer sequences give
\[
 0\longrightarrow E'(\mathbb Q)^{\wedge3}
 \longrightarrow H^1_{\rm cts}(G_{\mathbb Q},T_3E')
 \longrightarrow T_3H^1(\mathbb Q,E')\longrightarrow0.
\]
There is no degree-one inverse-limit error: the finite invariant
groups are Mittag--Leffler, and the quotients $E'(\mathbb Q)/3^n$
have surjective transition maps. A compatible sequence in the last
term has $\xi_n\in H^1(\mathbb Q,E')[3^n]$ and
$3\xi_{n+1}=\xi_n$. At every fixed $q\ne3$, local Tate duality
identifies the three-primary local Weil--Chatelet group with the
dual of the finite group $E'(\mathbb Q_q)^{\wedge3}$.
An infinitely compatible division sequence in this finite group is
zero. The real localization also has no three-primary part.

Global duality gives that the remaining localization $\xi_{n,3}$
annihilates all rational points. The precise applicable result is
\citet[I.6.26(b), p.~92]{CohMilneADT}, with all places and the
three-primary Tate--Shafarevich group, which E.111 proves zero.
By the local isomorphism just established, rational points are dense
in $E'(\mathbb Q_3)^{\wedge3}$. Local duality
\citep[I.3.3--3.5]{CohMilneADT} forces $\xi_{n,3}=0$.
Now every $\xi_n$ is everywhere locally trivial and hence zero by
E.111. This proves the absolute statement. The proof handles each
place of the complete compatible sequence and never assumes a
common finite set of ramified primes.
\end{proof}

\begin{result}{E.113}{Proved}{R6 C.3.4: exact necessary support tests}
For every original A.44 pair admitting a full geometrically
pro-three section, each fresh $A$-support prime $q_j$ and
$B$-support prime $r_j$ satisfies
\[
 3\nmid\operatorname{ord}_{E'(\mathbb F_{q_j})}(-2,1),\qquad
 3\nmid\operatorname{ord}_{E'(\mathbb F_{r_j})}(3,-6).
\]
These are necessary conditions for all original outputs. They do
not assert the existence of a section when the tests pass.
\end{result}
\begin{proof}
E.110 factors the section through $Y$; E.112 identifies its compact
class with $a\kappa(P_0)+b\kappa(T)$, $a\in\mathbb Z_3$,
$b\in\mathbb Z/3$. At a support prime $p>3$, A.31 identifies the
restriction of the full pro-three section with the ordinary local
point section, because one valuation is a nonzero integer.
A lift of a point section to a finite etale cover is a point section:
its fixed coset is a rational point of the fibre. Thus its proper
localization is the Kummer class of some $Z_p\in Y(\mathbb Q_p)$.

On $C$, the elementary map
$g_0(u,v)=(6uv,3(3u^3+2v^3))$ satisfies
$g_0=\psi\iota+P_0$, as exact substitution verifies. On $E'$ this
becomes $g_0\iota^{-1}\phi=[3]+P_0$. At an $A$-support prime,
the reduction has $u=0$, $-2v^3=1$, so $g_0=-T$; at a
$B$-support prime it has $v=0$, $3u^3=1$, so $g_0=T$.
The local three-primary Kummer classes consequently obey
\[
 (3a+1)\kappa_p(P_0)=-\kappa_p(T)\quad(A\text{-support}),\qquad
 (3a+1)\kappa_p(P_0)=\kappa_p(T)\quad(B\text{-support}).
\]
Because $3a+1$ is a three-adic unit, multiplication by three gives
$3\kappa_p(P_0)=0$. The displayed equations reduce to
$\kappa_p(P_0+T)=0$ or $\kappa_p(P_0-T)=0$.
Good reduction at $p>3$ identifies the three-primary completion
with the three-primary subgroup of $E'(\mathbb F_p)$.
Finally the group law gives $P_0+T=(-2,1)$ and $P_0-T=(3,-6)$.
This proves the tests and their exact quantifiers.
\end{proof}

\begin{result}{E.114}{Proved}{R6 C.3.5: infinitely supported pairs passing the compact tests}
The original A.44 recursion can be carried out with the additional
conditions that $3$ is a noncube modulo every $q_j$ and $2$ is a
noncube modulo every $r_j$. Every resulting pair passes all tests
of E.113. This confirms the stated compact-test survival theorem;
the statement that this proves full pro-three survival is unsupported.
\end{result}
\begin{proof}
At the $q_j$ choice, let $F$ be the cyclotomic field used for the
finite congruence conditions in E.102. It is abelian over $\mathbb Q$
and contains $\mu_3$. Rational cube classes inject into
$F^\times/F^{\times3}$: if a noncube rational number had a cube root
in an abelian extension, its nonnormal cubic root field would be a
subfield of an abelian Galois extension and hence normal, a contradiction.
The classes $3$ and $b_j=-2\prod_{i<j}r_i^{n_i}$ are independent,
since all $n_i$ are divisible by three and the valuations at three
and two distinguish them. The radical extension
$F(b_j^{1/n_j})/F$ is cyclic and has a unique cubic subextension,
$F(b_j^{1/3})$. Hence it is disjoint from $F(3^{1/3})$ over $F$.
The compatible automorphism used in E.102 can therefore be extended
with a prescribed nonzero action on $3^{1/3}$. Chebotarev supplies
infinitely many allowable $q_j$, retaining all old congruences and
the noncube condition. At the $r_j$ choice the same argument uses
the independent classes $2$ and $a_jq_j^{n_j}=3$ modulo cubes.
Here the radical splitting exponent is $n_{j+1}=(j+5)!$, while
the exponent with which $r_j$ enters $B$ remains $n_j=(j+4)!$.
This is the original recursion, not a changed exponent convention.

At an $A$-support prime the imposed radical splitting makes $-2$
a cube and $3$ is a noncube. At a $B$-support prime $3$ is a cube
and $2$ is a noncube. All these primes are $1\pmod3$.
The polynomial $3x(x^3+36)$ then shows
$E'(\mathbb F_p)[3]=\langle T\rangle$: $-36$ is a noncube under
either set of conditions. Thus the three-primary subgroup is cyclic.
The rational automorphism $\rho(x,y)=(\zeta_3x,y)$ satisfies
$1+\rho+\rho^2=0$ as an elliptic endomorphism. A cyclic group
whose order is divisible by nine cannot admit such an automorphism,
because $X^2+X+1$ has no root modulo nine. Consequently
$E'(\mathbb F_p)\{3\}=\langle T\rangle$ exactly.

The curves $E,E'$ are isogenous and have equal finite-field point
counts. For completeness, this equality can be obtained without a
point-count conjecture: both have the same rational Tate module for
any auxiliary prime other than $3,p$, so their Frobenius characteristic
polynomials and $\det(1-F)$ agree. The kernel of $\psi$ over
$\mathbb F_p$ has order three, hence $\psi E(\mathbb F_p)$ is the
prime-to-three subgroup of $E'(\mathbb F_p)$. Its Kummer test is
$Y-3$. At $P_0+T$ it is $-2$ and at $P_0-T$ it is $-9$.
The former is a cube at an $A$-support prime; the latter is a cube
at a $B$-support prime because $3$ is a cube and $-1$ always is.
The two points therefore have order prime to three as required.

This argument proves nonemptiness of the compact-test-passing
infinite-support family. It constructs neither a continuous section
of $\pi_1(Y)^{(3)}$ nor compatible sections of its nilpotent tower.
The exact finite examples in the accompanying program illustrate
only the finite-field test; they are not the recursive prime sequence.
\end{proof}

\subsection{The affine abelian quotient and the boundary norm calculation}

\begin{result}{E.115}{Proved}{R6 C.3.6--C.3.7: the exact generalized-Jacobian image}
Write $D=D_0\sqcup D_1\sqcup D_\infty=E'\setminus Y$,
$K_i=\Gamma(D_i,\mathcal O)$, and $J=\operatorname{Jac}(E',D)$.
The geometric abelianization of $\pi_1(Y)$ is $TJ$, with torus
\[
 \mathcal T=(\operatorname{Res}_{D/\mathbb Q}\mathbb G_m)/\mathbb G_m,
 \qquad 1\longrightarrow\mathcal T\longrightarrow J
       \longrightarrow E'\longrightarrow1.
\]
The exact image on the original three-primary abelian coordinates is
\[
 \begin{split}
 A&=3\left(u_0^a\,
        \frac{N_{K_0/\mathbb Q}\alpha_0}
             {N_{K_\infty/\mathbb Q}\alpha_\infty}\right)^3,\\
 B&=-2\left(v_0^a\,
        \frac{N_{K_1/\mathbb Q}\alpha_1}
             {N_{K_\infty/\mathbb Q}\alpha_\infty}\right)^3,
 \end{split}
\]
where $a\in\mathbb Z_3$, $\alpha_i\in(K_i^\times)^{\wedge3}$,
$u_0=253/655$, and $v_0=-488/655$. The full profinite affine
abelian image contains $(3C^{27},-2D^{27})$ for arbitrary completed
$C,D$. At the rational prime two the proper three-primary class
has no additional restriction: the three rational sheets over any
ordinary point with the required cube classes realize all classes
of $E'(\mathbb Q_2)^{\wedge3}=\langle T\rangle$.
\end{result}
\begin{proof}
For a smooth curve with reduced modulus, the degree-zero rigidified
Picard functor gives the stated torus exact sequence. The precise
representability and exact sequence are
\citet[Proposition 1.6, (1.4.1), p.~954]{AuditJordanRibetScholl}.
The pointed generalized Albanese map $j:Y\to J$ induces the
geometric abelianization; this is
\citet[Proposition A.8(iii), pp.~232--234]{AuditMochizuki2012}.
There is no discarded torsion here: the geometric fundamental group
of the affine genus-one curve with $27$ punctures is free profinite
of rank $2+27-1=28$, whose abelianization is torsion-free.

The functions $u,v$ have divisors $D_0-D_\infty$ and $D_1-D_\infty$,
respectively, and equal one at the origin. Indeed their cubes are
$t/3$ and $(1-t)/(-2)$, whose orders at the boundary are three.
Under the generalized Albanese universal property these functions
extend to homomorphisms $J\to\mathbb G_m$. On the torus their
characters are respectively the ratios of block norms from
$K_0,K_\infty$ and $K_1,K_\infty$. This also follows directly from
Weil reciprocity for a principal divisor and a rigidified line bundle.
The block-constant map $(s,r)\mapsto(s\text{ on }D_0,
r\text{ on }D_1,1\text{ on }D_\infty)$ has character values
$s^9,r^9$. Its effect on the original coordinates is thus
$s^{27},r^{27}$, proving the full profinite inclusion.

The torus Tate-module exact sequence gives
\[
 H^1(G_{\mathbb Q},T_3\mathcal T)
   =\left(\bigoplus_i(K_i^\times)^{\wedge3}\right)
        /(\mathbb Q^\times)^{\wedge3}.
\]
Here the only possible obstruction to the displayed surjectivity is
the kernel of restriction in $H^2(G_{\mathbb Q},\mathbb Z_3(1))$.
Kummer sequences identify that group with $T_3\operatorname{Br}(\mathbb Q)$:
the first-cohomology transition maps are surjective, so there is no
derived-limit term. This Tate module is torsion-free. Restriction
followed by the sum of corestrictions is multiplication by $27$,
so the kernel is zero. The first-cohomology diagonal injection is
also injective by restriction--corestriction and the torsion-freeness
of $(\mathbb Q^\times)^{\wedge3}$.

The classes of the rational points $P_0,T\in Y(\mathbb Q)$ lift
the generators of $H^1(G_{\mathbb Q},T_3E')$ from E.112.
Consequently every class in $H^1(G_{\mathbb Q},T_3J)$ is the sum
of a toric class and $a\kappa(j(P_0))+b\kappa(j(T))$, with
$a\in\mathbb Z_3$ and an integer representative $b$ modulo three.
There is no extra stabilizer from $H^0(G_{\mathbb Q},T_3E')$, which
is zero since the rational three-primary torsion is finite.
The exact map of E.110 gives $u(T)=v(T)=1$, and substitution at
$P_0$ gives $u_0,v_0$ as displayed; these identities were recomputed
in the retained program. Applying the two characters and cubing
now gives exactly the stated image. Choosing a different lift of
$b\pmod3$ only changes the toric class and does not change this image.

Finally E.111 already proves that $E(\mathbb Q_2)$ is uniquely
three-divisible and $\phi$ is surjective there. Its kernel is the
three rational points $O,\pm T$. Their distinct compact classes
exhaust $E'(\mathbb Q_2)^{\wedge3}$: multiplication by three on the
formal pro-two subgroup is an isomorphism, and the finite
three-primary quotient is precisely the kernel of $\phi$.
Equivalently, if $E'(\mathbb Q_2)$ had further three-primary
torsion, its image in the uniquely divisible compact group
$E(\mathbb Q_2)$ would give nonzero three-torsion there.
Every local cubic point therefore has three rational lifts differing
by $T$, and their compact classes give all three possibilities.
\end{proof}

\begin{result}{E.116}{Proved}{R6 C.3.8 and C.3.10: the boundary norms have no obstruction}
Each $K_i$ is a field of degree nine. Its normal closure has group
$G=H_3\rtimes C_2$, where the involution inverts the two generators
of $H_3$ and fixes its center; the point stabilizer is $H\simeq S_3$.
For $X=\operatorname{coker}(\mathbb Z\to\mathbb Z[G/H])$,
\[
 H^1(G,X)=H^2(G,X)=0.
\]
Each boundary field satisfies the Hasse norm principle, its norm-one
torus satisfies weak approximation, and
\[
 N:(K_i\otimes\mathbb Q_3)^\times\longrightarrow\mathbb Q_3^\times
 \quad\hbox{is surjective}.
\]
In particular, no additional congruence modulo $27$ is required in
the norm-lifting construction.
\end{result}
\begin{proof}
For the boundary over zero, the geometric coset set is
$H_3/\langle x\rangle$. The action twisted by the point-three
section has image $G$: over $\mathbb Q(\mu_3)$ the Kummer classes
$3,-2$ are independent, and surjectivity onto $H_3^{\rm ab}$ forces
surjectivity onto $H_3$, because its Frattini subgroup is its center.
The cyclotomic action supplies the involution. Its stabilizer is
$\langle x,\tau\rangle\simeq S_3$, with trivial core. Indeed the
intersection of the conjugates of its order-three subgroup is trivial,
and any nontrivial normal subgroup meeting its involution would have
a nontrivial commutator in that order-three intersection. Thus the
coset action is faithful and transitive of degree nine.
The transformations $t\mapsto1-t$ and $t\mapsto1/t$ give the other
two boundaries from the same zero-boundary calculation for point
values $-2$ and $1/3$. Their Kummer pairs are $(-2,3)$ and
$(1/3,2/3)$, again independent. This argument is applied to each
boundary field separately; no identification of their three normal
closures is asserted.

We prove the needed group cohomology. The presentation
\[
 H_3=\langle x,y\mid x^3=y^3=1,
 [x,[x,y]]=[y,[x,y]]=1\rangle
\]
has four relators. In a central extension by $\mathbb F_3$, chosen
lifts $X,Y$ yield four central relator values in $\mathbb F_3$.
Changing lifts does not change those values; if all four vanish,
the presentation supplies a splitting. Thus they give an injective
linear map $H^2(H_3,\mathbb F_3)\hookrightarrow\mathbb F_3^4$.
The extension has nilpotency class at most three. Replacing the
lifts by $X^{-1},Y^{-1}$ negates both cube values and both
weight-three commutator values. Therefore the involution acts by
$-1$ on $H^2(H_3,\mathbb F_3)$. Restriction identifies
$H^2(G,\mathbb F_3)$ with the invariant subspace, which is zero.
The coefficient exact sequence then gives $H^3(G,\mathbb Z)[3]=0$.
Its two-primary part restricts injectively to a Sylow $C_2$, since
the index is $27$, and $H^3(C_2,\mathbb Z)=0$. Finite-group
cohomology is killed by $|G|=54$, so $H^3(G,\mathbb Z)=0$.

The abelianizations of both $G$ and $H$ are $C_2$, and restriction
of the sign character is an isomorphism. Hence
\[
 H^2(G,\mathbb Z)=\operatorname{Hom}(G,\mathbb Q/\mathbb Z)
   \longrightarrow H^2(H,\mathbb Z)
\]
is an isomorphism. Shapiro's lemma and the exact sequence defining
$X$ now give both $H^1(G,X)=0$ and $H^2(G,X)=0$.
The precise global norm-theory comparison is
\citet[Theorems 2.1--2.3, (2.4), (2.6), pp.~496--498]{AuditMacedoNewton}:
the Hasse-norm obstruction is dual to a kernel in $H^2(G,X)$,
and the dual of the weak-approximation obstruction is a quotient
of the cyclic-restriction kernel in the same group, as in their
Proposition 2.4, (2.8). Both obstructions therefore vanish.

At three the same finite Galois image is still all of $G$.
The classes $3,-2$ are independent in
$\mathbb Q_3^\times/(\mathbb Q_3^\times)^3$: the first has valuation
one and the second is the noncube unit $7\pmod9$.
They remain independent after adjoining $\mu_3$, by
restriction--corestriction for the degree-two extension. The
transformed pairs for the other boundaries remain independent too.
Thus the local normal closure has group $G$, and each local degree-nine
algebra is a field with the same stabilizer $H$.
The action on $X$ factors through $G$. Inflation is an isomorphism
in degree one because a continuous homomorphism from the remaining
profinite kernel to the discrete torsion-free lattice $X$ is zero.
Therefore $H^1(\mathbb Q_3,X)=0$.
Local torus duality
\citep[I.2.4, pp.~28--29]{CohMilneADT} gives
$H^1(\mathbb Q_3,R^1_{K_i/\mathbb Q}\mathbb G_m)=0$.
Hilbert 90 in the norm exact sequence identifies this last group
with the cokernel of the asserted norm map, proving surjectivity.
\end{proof}

\begin{result}{E.117}{Proved}{R6 C.3.9--C.3.10: all compact-test-passing pairs lift to the affine abelian quotient}
Every original A.44 output satisfying the support tests in E.113
has a full profinite section of
$\pi_1(Y)/\overline{[K,K]}$, where
$K=\pi_1(Y_{\overline{\mathbb Q}})$, with the prescribed original
abelian pair. The proper elliptic component of the constructed
section is zero. This includes every pair constructed in E.114.
It does not construct a full pro-three section of $Y$.
\end{result}
\begin{proof}
The boundary identity in E.113 identifies the three etale degree-nine
divisors as
\[
 D_0=[3]^{-1}(-P_0-T),\qquad
 D_1=[3]^{-1}(-P_0+T),\qquad
 D_\infty=[3]^{-1}(-P_0).
\]
They are unramified outside two and three, because $E'$ has good
reduction elsewhere and $[3]$ is etale there. At an $A$-support
prime $q_j$, the support test says $-P_0-T$ has prime-to-three order
after reduction. It is therefore divisible by three over
$\mathbb F_{q_j}$. The corresponding point of $D_0$ lifts by
etaleness over $\mathbb Z_{q_j}$, so $K_0\otimes\mathbb Q_{q_j}$
has a degree-one factor and its norm is surjective. At every other
prime greater than three, $q_j$ is a unit and is a norm from the
unramified boundary algebra: norms on units of an unramified local
field extension are surjective. At two, every odd unit is a ninth
power, and the norm of its scalar ninth root is the required unit.
At three, use E.116. At the real place an odd-degree algebra has
a real factor, whose norm supplies any required sign.
Thus $q_j$ is everywhere a local norm from $K_0$, and E.116 gives
an actual $z_j\in K_0^\times$ with $N z_j=q_j$.
The same argument gives $w_j\in K_1^\times$ with $Nw_j=r_j$.

The products
\[
 \alpha_0=\prod_j z_j^{n_j/3},\qquad
 \alpha_1=\prod_j w_j^{n_j/3},\qquad\alpha_\infty=1
\]
converge in the respective three-power completions: at any fixed
coefficient exponent all sufficiently late factors are literally
trivial, since $v_3(n_j/3)\to\infty$. Their norm cubes are $A/3$
and $B/(-2)$. The toric class of E.115, with $a=0$, therefore gives
the required three-primary affine abelian section and zero proper
component. For every coefficient prime $\ell\ne3$, multiplication
by $27$ is invertible in the $\ell$-primary Kummer completion.
The block-constant torus map of E.115 supplies the prescribed
coordinates in that component, again with zero proper component.
Taking their product is a continuous cocycle for $TJ$, since maps
into a product are continuous exactly when all component maps are.
The rational origin splits the arithmetic extension, so this cocycle
is an actual section of the stated quotient. No nonabelian lifting
step has been inferred from the existence of this cocycle.
\end{proof}


\subsection{A verified failure within the compact-passing family}

\begin{result}{E.118}{Proved}{R6 C.3.11 and its existential premise C.2.7}
There is a recursively specified original A.44 pair in the
compact-test-passing family of E.114 whose unrestricted integral
class-three section fibre is empty. Every integral class-two section
has a nonzero third obstruction at the second $A$-support prime
$q_1$, for coefficient prime five, detected modulo $25$.
The full profinite affine abelian fibre over this pair is nonempty,
whereas the fibre for $\pi_1(Y)/\overline{[[K,K],[K,K]]}$ is empty.
This is an exclusion at coefficient prime five; it does not decide
whether the same pair has a full pro-three section.
\end{result}
\begin{proof}
We reconstruct the existential premise, rather than use the previous
coefficient-five assertion as an authority. Choose the first two
primes $q_0,r_0$ by E.114 and write
$a_1=3q_0^{24}$, $b_1=-2r_0^{24}$. At stage one the next exponent
is $120$. Choose its old local roots
$a_1\xi_p^{120}+b_1\eta_p^{120}=1$. Select rational integers
$P,R$, units at $q_0,r_0$, with
$P\equiv\xi_p$, $R\equiv\eta_p\pmod p$ for those two primes.
Fix integral five-adic Kummer cocycles $\alpha,\beta,p,r$ for
$a_1,b_1,P,R$ and form $a^*=\alpha+120p$, $b^*=\beta+120r$.
At $q_0$, $b^*=\delta u_0$ integrally, since
$b_1R^{120}$ is a principal unit. At $r_0$, similarly
$a^*=\delta v_0$. The primitives are unique, since the integral
weight-one invariant groups are zero.

The finite cup $\bar\alpha\cup\bar\beta$ is globally zero.
At each prime the original induction supplies a unit solution of
the exponent-$120$ diagonal equation, so its localization is the
Steinberg cup; the finite global detection proof is A.25.
The residue surjection of A.25 lets us choose a central primitive
$h_*$ with inertia values
\[
 h_*(T_{q_0})=24\bar u_0,\qquad
 h_*(T_{r_0})=-24\bar v_0,
\]
and with zero central inertia at all other primes outside
$\{2,3,5\}$. Only finitely many adjustments are required, because
a continuous finite-valued cochain factors through a finite quotient.
The unramified uniqueness group is zero:
\[
 H^1(G_{\mathbb Q,\{5,\infty\}},\mathbb F_5(2))=0.
\]
Indeed the integral first group is zero by A.26, and the integral
second group is $K_2(\mathbb Z)\{5\}=0$ by A.25; the coefficient
exact sequence gives the assertion. At two and three there is no
extra residue because $p^2-p$ is a unit modulo five. Thus $h_*$
determines, up to central conjugacy, a finite homomorphism
\[
 \rho_*:G_{\mathbb Q}\longrightarrow H_5\rtimes\mathbb F_5^\times.
\]
Its image is the full group. Over $\mathbb Q(\mu_5)$ the classes
$a_1,b_1$ are independent by their valuations at three and two;
the Heisenberg commutator then forces the full $H_5$ image.
The semidirect product has commutator subgroup $H_5$, since a
cyclotomic scalar different from one generates the two abelian
weight-one coordinates by commutators. Consequently its fixed
field $L_*$ has maximal abelian subfield $\mathbb Q(\mu_5)$.

This finite class is fixed before $q_1,r_1$ are chosen. To prove
this, require all later $q_j,r_j$, $j\ge2$, to be $1\pmod{q_0r_0}$.
At either old prime the differences of the completed Kummer
cocycles from $a^*,b^*$ are then
$5\delta e_p$, $5\delta f_p$ with integral primitives.
In fact $q_1/P,r_1/R$ and every later prime are principal units
there, and the primitive series converges because $n_j/5\to0$
five-adically. The actual tame relation makes the inertia of any
integral class-two lift reduce to the two prescribed values.
Every later support valuation is divisible by five, so its central
inertia is zero modulo five; at zero-valuation primes it is zero
integrally. The difference from $h_*$ is therefore in the zero
unramified uniqueness group. This proves the exact prefix assertion
R6 C.1.20 with integral primitives retained before reduction.

We next choose $q_1$ with a central Frobenius defect while preserving
the noncube condition. At the old prime five choose $\xi_5=1$.
The original invariant has
$v_5(a_1+b_1-1)\ge6$ and $v_5(120)=1$, so the root for $\eta_5$
can be chosen $1\pmod{5^5}$. It satisfies the original stronger
next-stage congruence as well. The original stage-one CRT class
$Q_0$ therefore satisfies $Q_0\equiv1\pmod{720}$ and
$Q_0\equiv1\pmod{25}$. Let $F=\mathbb Q(\mu_M)$ be its
cyclotomic field and set
\[
 R_b=F(b_1^{1/120}),\qquad
 R_*=F(b_1^{1/120},a_1^{1/5}),\qquad
 C_*=F(3^{1/3}).
\]
The original radical compatibility proof E.102 gives an automorphism
of $R_b$ extending the specified CRT action on $F$ and fixing
$b_1^{1/120}$. The rational classes $a_1,b_1$ remain independent
modulo fifth powers in $F$. For a nontrivial rational fifth-power
class has a radical extension of degree five which is nonnormal;
it cannot lie in an abelian extension of $\mathbb Q$.
The cyclic extension $R_b/F$ has unique degree-five subfield
$F(b_1^{1/5})$, so it is disjoint from $F(a_1^{1/5})$.
Extend the chosen automorphism to $R_*$ fixing this fifth root too.

Since $L_*\cap F=\mathbb Q(\mu_5)$, we have
$\operatorname{Gal}(L_*F/F)=H_5$. The extension $R_*/F$ is abelian,
so its intersection with $L_*F$ is exactly the maximal abelian
subextension
$F(a_1^{1/5},b_1^{1/5})$, which is already contained in $R_*$.
Thus $[L_*R_*:R_*]=5$. A chosen nonzero central element of $H_5$
agrees with the radical-fixing automorphism on this intersection,
so both actions extend to $L_*R_*$. Finally $C_*$ is disjoint
from $R_*$: the classes $3,b_1$ are independent modulo cubes and
the extra fifth-root extension has degree prime to three.
It remains disjoint after adjoining $L_*$, since the added degree
is five. We may therefore impose a nontrivial cubic action on
$3^{1/3}$ simultaneously with the nonzero central action on $L_*$.

Chebotarev, in the exact form used in E.102, now supplies $q_1$
beyond all original bounds, with the original CRT and radical
conditions and also
\[
 q_1\equiv1\pmod{25},\qquad
 3\notin(\mathbb F_{q_1}^\times)^3,\qquad
 \rho_*(F_{q_1})=(0,0,c;1),\quad c\ne0\text{ in }\mathbb F_5.
\]
These are conjugacy-invariant conditions. Choose $r_1$ and all
later primes by E.114. The congruences $q_j,r_j\equiv1\pmod{q_0r_0}$
for $j\ge2$ can be imposed through the old roots: after the first
continuation stage the coefficient sum is one modulo each old
prime; simple Hensel lifting applies when it does not divide $n_j$,
and the original inequality $H_j(p)>2v_p(n_j)$ applies otherwise.
Both give near-one roots, compatible with the stronger next-stage
condition. Hence this is a legal infinite recursion, and every
support prime satisfies the compact tests. Selecting least eligible
primes makes it recursive. The required marked finite extension
$L_*$ can itself be specified by enumerating finite Galois realizations
and checking the finite ramification and inertia data; existence
and uniqueness of its kernel have just been proved. The argument
does not pretend to print its enormous prime values.

For any integral class-two section of this pair, the prefix
calculation gives at $q=q_1$ the integral Frobenius coordinates
$\alpha\equiv0\pmod5$ and $\gamma\equiv c\ne0\pmod5$.
Central coboundaries do not change the latter residue since
$q\equiv1\pmod5$. The ordinary local data are
$m=120$, $b=\delta u$, and $h(T_q)=120u$.
The actual third residue from B.23 or B.102 is
\[
 Z_{D,q}=120(\gamma-q\alpha u)
          +(q^2-1)u\binom{120}{2}.
\]
The first term has five-adic valuation exactly one, while the
second is divisible by $25$. Hence $v_5(Z_{D,q})=1$, whereas
$v_5(q^2-1)\ge2$. The local third obstruction is nonzero, already
with coefficients modulo $25$. Every integral class-two class is
covered by this argument; their existence is A.25. Thus the first
empty unrestricted integral nilpotent fibre is class three.

E.117 supplies a section of the full profinite affine abelian
quotient for the same pair. Write $K'=\overline{[K,K]}$ and
$K''=\overline{[K',K']}$. For the full geometric tripod group
$\Pi$, the actual cover inclusion gives
$K'\subseteq\Gamma_2\Pi$ and
$K''\subseteq[\Gamma_2\Pi,\Gamma_2\Pi]\subseteq\Gamma_4\Pi$,
with closures understood. Therefore an arithmetic section of
$\pi_1(Y)/K''$ would push to an integral class-three section of
$U$, which has just been excluded. This proves the asserted empty
metabelian fibre and confirms R6 C.3.11. The proof uses a
five-primary obstruction and supplies no three-primary obstruction;
replacing its full profinite quotient by a pro-three quotient would
discard precisely the witness used here.
\end{proof}

\begin{result}{E.119}{Proved}{The selected local-fibre exclusion in R6 C.3.1 and C.2.6}
Original A.44 primes can also be chosen so that the ordinary local
point at $q_0$ has no $\mathbb Q_{q_0}$-rational lift to $Y$.
For these selected outputs the unrestricted integral class-three
fibre is empty, with the obstruction at coefficient prime three.
Consequently no full geometrically pro-three section exists over
these outputs. These selected outputs are outside the compact-passing
family of E.114.
\end{result}
\begin{proof}
Let $L_3$ be the finite field cut out by the point-three section in
$H_3\rtimes\mathbb F_3^\times$. Independence of $3,-2$ over
$\mathbb Q(\mu_3)$ gives full image, and the commutator calculation
used in E.116 gives maximal abelian subfield $\mathbb Q(\mu_3)$.
At the initial A.44 stage all old roots may be chosen equal to one,
because $3-2=1$ identically. Let $F=\mathbb Q(\mu_M)$ be the
resulting cyclotomic field. Its allowed Frobenius is the identity,
and the original modulus is divisible by nine. Set
$R=F((-2)^{1/24},3^{1/3})$. It is abelian over $F$, and
\[
 L_3F\cap R=F((-2)^{1/3},3^{1/3}),\qquad [L_3R:R]=3.
\]
Indeed the left intersection is contained in the maximal abelian
subextension of the Heisenberg extension $L_3F/F$, and the
displayed Kummer field is already contained in $R$.
A nonzero central automorphism of $L_3F/F$ therefore extends to
the identity on $R$. Chebotarev supplies an eligible $q_0$, beyond
all original bounds, which is $1\pmod9$, splits the required
radicals, and has nonzero central Frobenius in $L_3$.
Continue the original legal recursion; no further restrictions are
needed for this witness.

The ordinary local point has valuation $24$ at $q_0$. Both its
modulo-three Kummer coordinates are zero, because $3,-2$ are cubes
there and every new exponent is divisible by three.
A.31 identifies its full local pro-three section with the cyclic
boundary section. Its image in $H_3$ is therefore trivial. The
point-three section, by construction, has nonzero central Frobenius;
the local cyclotomic action on $H_3$ is trivial. These two quotient
sections cannot be geometrically conjugate. A rational point of the
cover fibre is exactly a fixed coset, or equivalently a conjugacy
between these sections. The fibre consequently has no rational point.
E.110 and A.31 then exclude any full pro-three section with this pair.

One can already see the failure at integral class three. The
modulo-three class-two possibilities are $h_3+k\eta$. Those with
$k\ne0$ fail at the rational prime two by E.110. For $k=0$ the
Frobenius data at $q_0$ have $\alpha\equiv0\pmod3$ and
$\gamma\not\equiv0\pmod3$. The first actual third residue is
\[
 24(\gamma-q_0\alpha u)+(q_0^2-1)u\binom{24}{2}.
\]
Its three-adic valuation is one, whereas
$v_3(q_0^2-1)\ge2$. It is nonzero already modulo nine. Every
integral class-two section is therefore excluded, while such
sections do exist by A.25. This proves the precise first-failure
claim and the selected local-fibre portion of R6 C.3.1.
The compact-passing construction imposes that $3$ be a noncube at
every $A$-support prime, so these choices cannot lie in that family.
\end{proof}

\begin{result}{E.120}{Open}{The exact pro-three question left by the elliptic audit}
It remains undetermined whether any original A.44 pair constructed
in E.114 admits a continuous section of the full geometrically
pro-three arithmetic fundamental group of $U$. Equivalently, after
E.110, the question is whether some affine abelian class in the exact
image of E.115 lifts through every compatible nonabelian quotient
of the actual pro-three group of $Y$.
\end{result}
The open step is not a Selmer-rank calculation or a boundary norm
principle: E.111--E.117 prove those statements completely. Nor is
the coefficient-five failure in E.118 an obstruction in a pro-three
group. The stronger proper localization of E.112 supplies a genuine
necessary gate: prescribing the proper class at $\mathbb Q_3$
fixes its global coefficients $(a,b)$ uniquely. It supplies no
local point theorem for the punctured curve at the coefficient
prime three. A.31, which proves local point matching at a support
prime different from the coefficient prime, has no such assertion
at $q=\ell=3$.

At all support primes $q\ne3$, the stronger integral class-two
matching condition of B.24 is necessary for full pro-three
survival. It compares the canonical global lift to the actual
ordinary local point before multiplying by the support valuation;
the compact tests only see its proper elliptic image after reduction.
One must still construct a legal prime sequence satisfying those
matching conditions, or prove a necessary failure for every legal
sequence, and then control higher compatible obstructions. The
Chebotarev construction of E.114 controls cubic residue symbols and
the finite three-primary reduction of the proper curve. It does not
prove that additional simultaneous integral matching assertion.

\begin{result}{E.121}{Proved}{A new three-primary exclusion inside the compact-passing family}
Let $\mathscr P$ be the original A.44 family with $3$ a noncube
at every $A$-support prime and $2$ a noncube at every $B$-support
prime. Let $\mathscr A$ be the entire original A.44 family and
$\mathscr A_c$ its compact-test-passing subspace. Every legal
finite prefix in any of $\mathscr A,\mathscr A_c,\mathscr P$ has an extension
whose geometrically pro-three class-three section fibre is empty.
For prefixes in $\mathscr A_c$ or $\mathscr P$, the extension
still passes all compact elliptic tests and has the full profinite
affine abelian section of E.117. A prefix in $\mathscr P$ extends
within $\mathscr P$.

More precisely, at a new stage $j$ put $n=n_j$ and $f=v_3(n)$.
One can choose its $A$-support prime $q=q_j$ with
$e=v_3(q-1)>f$. Among three explicitly related choices for the
local $n$th root specifying the subsequent $r_j$ CRT condition,
at least two force a nonzero third obstruction modulo $3^{f+1}$
for the unique class-two reduction that can pass the test at two.
The other two global pro-three class-two classes already fail at
two. Thus examples with a local three-primary obstruction are
dense in each of these three spaces. The locally excluded set is
also open in the legal-prefix topology; in each space its complement
is a closed nowhere dense necessary remainder for full pro-three
survival. This proves the density conclusions of R6 C.2.8--C.2.9
by a three-primary argument, without certifying their distinct
claim that a coefficient prime at least five may always be used.
\end{result}
\begin{proof}
We keep literal global Kummer cocycles of the form
\[
 a=a_3+\sum_i n_i\kappa(q_i),\qquad
 b=b_{-2}+\sum_i n_i\kappa(r_i),
\]
where $a_3,b_{-2}$ are fixed representatives for the point-three
section. These sums converge uniformly. Because every $n_i$ is
divisible by three, the literal reductions are $a_3,b_{-2}$.
The inertia argument of E.110 proves that every integral class-two
lift reduces to $h_3+k\eta\pmod3$ up to a central coboundary.
The local third test at two forces $k=0$. At a prime $q\equiv1\pmod3$
the central coboundary vanishes on Frobenius modulo three. Thus
for a lift which could reach class three,
\[
 \alpha=a(F_q)\equiv a_3(F_q),\qquad
 \gamma=h(F_q)\equiv h_3(F_q)\pmod3.
\]
These two residues are fixed before $r_j$ is selected. Since $3$
is a noncube at $q$, $\bar\alpha\ne0$.

To choose $q$ with $e>f$ without altering the prescribed prefix,
set the old root $\xi_3=1$ at this new stage. The coefficient
$b_j$ is a three-adic unit, and the original invariant
$v_3(a_j+b_j-1)\ge H_j(3)>2v_3(n_j)$ gives a solution
$\eta_3$ near one to
$a_j+b_j\eta_3^{n_j}=1$. Its precision also satisfies the original
next-stage root requirement. Therefore the allowed $q$ CRT class
is $1\pmod{3^{L_3}}$, where
$L_3=H_{j+1}(3)=2v_3(n_{j+2})+j+4>f$.
E.114's independent cubic Kummer extension retains the noncube
condition for $3$ in this CRT class. Choose any resulting eligible
prime $q$. In particular $q\nmid n$ and $\mu_{3^f}\subset\mathbb Q_q$.

The original construction chooses a unit $z\in\mathbb Z_q$ with
\[
 b_j z^n=1-a_jq^n
\]
and specifies only $r_j\equiv z\pmod q$ in its new-prime CRT
condition. Let $\zeta\in\mathbb Q_q$ be a primitive $3^f$th root
of unity. Each $z\zeta^k$, $k=0,1,2$, is an equally legal root.
All three corresponding CRT classes admit the original radical
splitting and the noncube condition for $2$, by the independent
field argument of E.114. Thus this is genuine freedom in the
original recursion, not a change to its exponent or splitting rule.

We compute exactly what this freedom changes. At $q$, after the
choice of $r=r_j$, the number $b_jr^n$ is a principal unit. Hence
its literal three-adic Kummer cocycle is $\delta u_r$ for a unique
integral $u_r\in\mathbb Z_3(1)$. If two allowable choices $r,r'$
have residue quotient $\zeta$, then
\[
 (q-1)(u_{r'}-u_r)
       =n\,\bigl(\kappa(r')-\kappa(r)\bigr)(F_q).
\]
The unit Kummer isomorphism at $q\ne3$ identifies the class of
the right-hand cocycle value before multiplication by $n$ with
an element of exact order $3^f$ in $\mathbb Z_3/(q-1)$.
It therefore has valuation $e-f$ modulo $3^e$.
Multiplication by $n$ and division by $q-1$ give a unit modulo
three. Changing a Kummer representative changes the quotient
by $ns$, $s\in\mathbb Z_3$, which is zero modulo three.
Thus the three roots $z,z\zeta,z\zeta^2$ make $\bar u_r$
range through all of $\mathbb F_3$; the conclusion does not depend
on which eligible prime in each fixed CRT class is chosen.

Since $\bar\alpha\ne0$ and $\bar\gamma$ is fixed, exactly one
of these three primitive residues satisfies
$\bar\gamma-\bar q\bar\alpha\bar u_r=0$.
Choose either of the other two root classes. At all subsequent
stages require $q_i,r_i\equiv1\pmod q$. This is available by the
near-one old-root choices: immediately after the present stage
$a_{j+1}+b_{j+1}\equiv1\pmod q$; while $q\nmid n_i$ use
simple Hensel lifting, and once $q\mid n_i$ the original strong
invariant supplies the required higher-precision root. The cubic
nonresidue conditions remain compatible by E.114 at every stage.

Each later $r_i$ is a principal $q$-unit, so its integral local
Kummer cocycle is $\delta w_i$ for some integral $w_i$.
The completed primitive is
$u=u_r+\sum_{i>j}n_iw_i$. The series converges and every summand
is divisible by three. Hence $\bar u=\bar u_r$, and the chosen
nonzero mismatch persists in the completed ordinary pair.
At this support prime the actual residue is
\[
 Z_{D,q}=n(\gamma-q\alpha u)
               +(q^2-1)u\binom n2.
\]
Its first term has valuation exactly $f$, while the second has
valuation at least $e+f$ because $n-1$ and two are three-adic
units. Thus $v_3(Z_{D,q})=f<e=v_3(q^2-1)$, and its obstruction
is nonzero already modulo $3^{f+1}$. The reductions $k\ne0$
were excluded at two. No pro-three class-two class can therefore
lift to class three, and a fortiori there is no full pro-three
section above this pair.

The construction uses no noncube property of the preceding primes:
their exponents alone give the common cube classes $(3,-2)$.
It therefore works after every legal finite prefix in $\mathscr A$.
If that prefix belongs to $\mathscr A_c$, all its already fixed
support primes keep their compact tests and the new noncube
conditions ensure the tests for all future primes. This proves
density in $\mathscr A_c$ as well as in $\mathscr P$.
For openness, a nonzero local continuous obstruction is
detected at a finite coefficient exponent, as here explicitly.
B.20 fixes the entire finite set of integral pro-three class-two
reductions on a sufficiently long prefix. The actual class-three
defect is an integral polynomial in these coordinates at three;
thus the exclusions at two and $q$ persist for every legal
continuation of such a prefix. This is the coefficient-three
version of the closedness argument in B.21. The complement of
this open dense excluded set is the stated closed nowhere dense
necessary remainder. Nothing here asserts that this remainder is
empty or that its points lift beyond class three.
\end{proof}

\begin{result}{E.122}{Proved}{Exact local matching at three in the affine abelian quotient does not force finite support}
There are ordinary original A.44 ghosts in $\mathscr P$ with
\[
 A_{\mathbb Q_3}=\kappa(3),\qquad
 B_{\mathbb Q_3}=\kappa(-2),
\]
whose full profinite affine abelian cover fibre has a continuous
section with zero proper elliptic class and whose restriction to
$G_{\mathbb Q_3}$ is the actual origin point section of
$Y\subset E'$, up to geometric conjugacy. Thus exact local matching
at the coefficient prime three, even in this full affine abelian
quotient, together with the proper rigidity of E.112 and all compact
support tests does not imply finite support.
Such examples may simultaneously be chosen to have empty pro-three
class-three fibre by E.121. This is a counterexample to that proposed
rigidity inference, not a full nonabelian ghost section.
\end{result}
\begin{proof}
We first impose the exact ordinary point at three while retaining
the original recursion. Write $f_j=v_3(n_j)$ and
$L_j=H_{j+1}(3)=2f_{j+2}+j+4$. In addition to the original
inductive conditions, maintain
\[
 \min\{v_3(a_j/3-1),v_3(b_j/(-2)-1)\}
       \ge f_j+f_{j+1}+1.
\]
This holds initially with infinite valuation. Both quotients lie
in $1+3\mathbb Z_3$. The three-adic logarithm gives unique roots
in that group
\[
 \xi_3=(3/a_j)^{1/n_j},\qquad
 \eta_3=((-2)/b_j)^{1/n_j}.
\]
They satisfy $a_j\xi_3^{n_j}+b_j\eta_3^{n_j}=1$ exactly, and
their differences from one have valuation at least $f_{j+1}+1$.
In particular they meet the original requirement that both old
roots be $1\pmod{3^{v_3(n_{j+1})}}$.
Choose the original CRT primes with these residues modulo $3^{L_j}$.
Then
\[
 \begin{split}
 a_{j+1}/3&=(q_j/\xi_3)^{n_j},\\
 b_{j+1}/(-2)&=(r_j/\eta_3)^{n_j},
 \end{split}
\qquad
 \min\{v_3(a_{j+1}/3-1),v_3(b_{j+1}/(-2)-1)\}
       \ge L_j+f_j.
\]
The last inequality is the logarithmic power estimate in
$1+3\mathbb Z_3$. It implies the required bound at the next
stage, since
\[
 L_j+f_j-(f_{j+1}+f_{j+2}+1)
      =f_{j+2}+f_j-f_{j+1}+j+3>0.
\]
It also preserves the original sum invariant at three; in fact
the error just obtained is at least $H_{j+1}(3)$. All other old
roots and conditions remain those of E.102. E.114 permits the
two noncube conditions with every resulting legal CRT class.
The separate errors tend to infinity, so the ordinary local limits
are exactly $3$ and $-2$. The support primes remain pairwise distinct,
so the completed pair remains a ghost.

Let $\alpha_0\in(K_0^\times)^{\wedge3}$ and
$\alpha_1\in(K_1^\times)^{\wedge3}$ be the norm elements from
E.117. Their norm cubes are $A/3$ and $B/(-2)$. Their local norms
at three are one: those cubes are one, and
$(\mathbb Q_3^\times)^{\wedge3}$ has no three-torsion.
Put $L_i=K_i\otimes\mathbb Q_3$, which E.116 proves is a field,
and let $T_i=R^1_{K_i/\mathbb Q}\mathbb G_m$.
Local norm surjectivity gives
\[
 1\longrightarrow T_i(\mathbb Q_3)
 \longrightarrow L_i^\times
 \xrightarrow{N}\mathbb Q_3^\times\longrightarrow1.
\]
The finite three-power quotient sequences are exact on the left
as well: if a power in the kernel has a $3^m$th root in
$L_i^\times$, the norm of that root is a $3^m$th root of unity
in $\mathbb Q_3^\times$, hence is one. All local finite power
quotients are finite, so their inverse limits give
\[
 \ker\bigl((L_i^\times)^{\wedge3}
          \longrightarrow(\mathbb Q_3^\times)^{\wedge3}\bigr)
           =T_i(\mathbb Q_3)^{\wedge3}.
\]
The local image of $\alpha_i$ therefore belongs to this kernel.

We verify the compatible approximation needed to remove it.
For an abelian group $M$ with dense image in the finitely generated
pro-three group $L$, the induced map $M^{\wedge3}\to L$ is
surjective. Given $x\in L$, choose $m_0\in M$ whose image equals
$x$ modulo $3L$. Choose a local third root of the difference,
approximate that root by an element of $M$ modulo $3L$, and
continue. Additively, this constructs $m_j\in M$ with
\[
 x-\sum_{j=0}^{N}3^j m_j\in3^{N+1}L.
\]
The series is Cauchy in every global quotient $M/3^mM$ because
all terms with $j\ge m$ vanish there. Its completed sum maps to
$x$. This proves surjectivity with compatible corrections, not
merely surjectivity onto separate local finite quotients.

Apply this argument with $M=T_i(\mathbb Q)$ and
$L=T_i(\mathbb Q_3)^{\wedge3}$. Weak approximation in E.116
gives the needed density; local torus structure gives that $L$
is a finitely generated pro-three group. We obtain a completed
global norm-one element $\lambda_i$ whose local image is
$\alpha_i^{-1}$. Its image in $(K_i^\times)^{\wedge3}$ has
global norm one, because each term of its defining compatible
product has norm one. Replacing $\alpha_i$ by $\alpha_i\lambda_i$
therefore preserves both global norms and makes each local image
at three equal to one. With $\alpha_\infty=1$, the resulting
toric three-primary cocycle has trivial localization at three
and still gives the prescribed pair with proper class zero.

For every coefficient prime $\ell\ne3$, use the block-constant
torus construction of E.117. Since the original pair localizes
exactly to $(3,-2)$ at three, the required global $27$th roots
localize to one there; multiplication by $27$ is an automorphism
in these coefficient components. Their toric localizations are
therefore zero too. The product of the component cocycles yields
a continuous full profinite affine abelian section with the
claimed zero localization. Relative to the origin splitting, zero
local cohomology class is exactly geometric conjugacy to its point
section.

Finally the above choice of $\xi_3$ already gives
$q_j\equiv1\pmod{3^{f_{j+1}+1}}$, so
$v_3(q_j-1)>f_j$. E.121 needs only this inequality, not the
literal equality $\xi_3=1$. Its choice among the roots
$z_q\zeta^k$ at the distinct new prime $q_j$ is independent of
the two chosen old roots at three. The later congruences modulo
$q_j$ preserving the mismatch can likewise be imposed together
with the exact-target roots at three. Hence the examples may be
chosen to fail pro-three class three while satisfying every
assertion already proved here. This isolates the missing genuinely
nonabelian matching requirement at the other support primes.
\end{proof}

\subsection{A nonempty boundary-matched remainder at coefficient three}

\begin{result}{E.123}{Proved}{Ordinary ghosts with all support boundary matches in class two}
There are pairwise distinct primes $q_j,r_j>3$ and integers
$n_j=(j+7)!$ such that
\[
 A=3\prod_{j\geq0}q_j^{n_j},\qquad
 B=-2\prod_{j\geq0}r_j^{n_j}
\]
are an ordinary-local completed unit pair with both coordinates of
infinite support, and
\[
 v_3(q_j^2-1)=1,\qquad v_3(r_j^2-1)=2
 \quad(j\geq0).
\]
The pair has a global integral pro-three class-two section which,
at every support place different from three, is the class-two image
of the actual local point section. This includes the old support
place two. Consequently it has a global integral pro-three
class-three section.

This is a new recursion, with fixed-root conditions in place of
complete radical splitting. It is not an output of the original
A.44 recursion, and it does not assert a section at every depth.
It proves that the coefficient-three boundary-matching gate of
B.24 has a nonempty infinitely supported remainder.
\end{result}
\begin{proof}
We give the finite field conditions, prove their simultaneous
realizability, and then verify the continuous lifting assertion.

\emph{The fixed point-three field.}
Let $H_9$ be the class-two Heisenberg group with both abelian
coordinates and its central coordinate reduced modulo nine.
Reduce the actual integral class-two cocycle of the rational point
$3$ modulo nine, and retain the cyclotomic character modulo nine.
Its finite splitting field $L_9$ has Galois group
\[
 G_9=H_9\rtimes(\mathbb Z/9)^\times.
 \tag{E.123.1}
\]
Here the multiplier $u$ acts on the two generators by their
$u$th powers and on the center by $u^2$. To verify the full image,
work first over $\mathbb Q(\mu_9)$. The two abelian characters
are the Kummer classes of $3$ and $-2$. The valuation at two,
where the cyclotomic extension is unramified, makes the $-2$
coordinate independent. Moreover $3$ is not a cube in any abelian
extension of $\mathbb Q$: such an extension containing a cube root
also contains the real cube root, whose degree-three field is not
normal. Its Kummer class has order nine in a field containing
$\mu_9$. Thus the two characters give $(\mathbb Z/9)^2$.
A subgroup of $H_9$ surjecting onto this abelian quotient contains
a central generator, by the commutator of two basis lifts, and
therefore is all of $H_9$. The full cyclotomic image then proves
(E.123.1).

The maximal abelian subfield of $L_9$ is $\mathbb Q(\mu_9)$.
Indeed the central subgroup is generated by commutators of the
two geometric generators, and a multiplier $-1$ kills both
geometric generators in the abelianization. Consequently, for
every cyclotomic field $F$ containing $\mu_9$,
\[
 L_9\cap F=\mathbb Q(\mu_9),\qquad
 \operatorname{Gal}(L_9F/F)=H_9.
 \tag{E.123.2}
\]
This verifies the needed independence at three without assigning
valuation one to $3$ in $\mathbb Q(\mu_9)$.

\emph{Fixed radical roots and the exact intersection.}
The radical-field proof in E.102 applies to every $24\mid h$
and $c=3w^8$ or $-2w^8$ with $(w,6)=1$. For
$K=\mathbb Q(\mu_h)$ and $L=K(c^{1/h})$ it proves
$L^{\rm ab}=K$ and $[L:K]=h/2$. If $F$ is a cyclotomic field
containing $K$, the same degree holds over $F$. Choose a root
$z^h=c$, so $z^{h/2}=\sqrt c$. A specified automorphism
$\sigma$ of $F/\mathbb Q$ extends to an automorphism of $F(z)$
fixing $z$ if and only if it fixes $\sqrt c$.
For sufficiency, use the irreducible polynomial
$X^{h/2}-\sqrt c$ over $F$; its coefficients are fixed by
$\sigma$, so $\sigma(z)=z$ defines the extension. Necessity is
immediate. Chebotarev applied to this automorphism supplies primes
in the specified cyclotomic class for which $c$ is an $h$th
power in the residue field. Complete splitting is unnecessary.

All exponents in the present recursion are divisible by nine.
If $c=b=-2\prod r_i^{n_i}$ and $9\mid h$, then $F(c^{1/h})$
contains $F((-2)^{1/9})$. Its intersection with $L_9F$ is
exactly this field:
\[
 L_9F\cap F(b^{1/h})=F((-2)^{1/9}).
 \tag{E.123.3}
\]
In fact the intersection is cyclic over $F$, because the radical
extension is cyclic. Every cyclic quotient of $H_9$ factors
through its exponent-nine abelianization, so its degree is at
most nine. The displayed ninth-root subfield already has degree
nine. The same argument gives
\[
 L_9F\cap F(c^{1/h})=F(3^{1/9})
 \quad\hbox{if }c=3\prod q_i^{n_i}.
 \tag{E.123.4}
\]
The degree assertions also follow directly from the preceding
Kummer independence argument over $F$: ramification at two in
a cyclotomic field has power-of-two index, and a pure cube root
of $3$ cannot belong to an abelian extension.

\emph{The modified local roots.}
Keep the invariants of E.102, but put
\[
 n_j=(j+7)!,\qquad
 H_j(p)=2v_p(n_{j+1})+j+3.
\]
At stage $j$ write $n=n_j$, $m=n_{j+1}$ and $a=a_j$, $b=b_j$.
Thus $a=3u^8$, $b=-2v^8$, the supports are disjoint, the old
diagonal curve $aX^n+bY^n=1$ is locally soluble with unit
coordinates, and the sum error has valuation at least $H_j(p)$
for $p\leq B(n)$. Start at $(a_0,b_0)=(3,-2)$.
Use the same finite set
\[
 S=\{p:p\leq B(m)\}\cup\{p:p\mid ab\},
 \qquad L_p=H_{j+1}(p),\qquad
 M=\prod_{p\in S}p^{L_p}.
\]
At every old place except two and three choose the old unit
solution as in E.102; at primes dividing $m$ it may still be
chosen near one. At two instead set $\eta_2=-1$ and choose
the root $\xi_2\in1+4\mathbb Z_2$ of
\[
 \xi_2^n=(1-b)/a.
\]
This is possible since the right side differs from one by an
element of valuation at least $H_j(2)$, strictly larger than
$v_2(n)+2$. The logarithm or the elementary root estimate gives
$v_2(\xi_2-1)\geq H_j(2)-v_2(n)>3$.

At three set $\xi_3=2$. Write $f=v_3(n)\geq2$ and
$d=a+b-1$. Since $n$ is even,
\[
 v_3(2^n-1)=f+1,
 \qquad
 v_3\left(\frac{1-a2^n}{b}-1\right)=f+2.
 \tag{E.123.5}
\]
For the second equality use $v_3(a)=1$, $v_3(b)=0$, and
$v_3(d)\geq H_j(3)>f+2$, so no cancellation is possible.
There is a unique $Y\in1+3\mathbb Z_3$ with
$Y^n=(1-a2^n)/b$; the logarithm gives $v_3(Y-1)=2$.
Set $\eta_3=-Y$. Because $n$ is even these roots still satisfy
the exact old unit equation, and
\[
 v_3(\xi_3+1)=1,\qquad v_3(\eta_3+1)=2.
 \tag{E.123.6}
\]

\emph{Choosing the new $A$ prime.}
Let the CRT class modulo $M$ be that of the old roots $\xi_p$.
It is $1\pmod8$ and $2\pmod9$. In particular its action fixes
$\sqrt{-2}$, hence $\sqrt b$. Put $F=\mathbb Q(\mu_M)$.
Prescribe on $F(b^{1/n})$ the extension of this cyclotomic
automorphism which fixes a chosen $n$th root of $b$.
On $L_9F$ prescribe, in the point-three Hall coordinates,
\[
 (a_0(F_q),b_0(F_q),h_0(F_q))=(0,0,0)\pmod9.
 \tag{E.123.7}
\]
The multiplier is the already specified CRT class modulo nine.
The full group (E.123.1)--(E.123.2) permits this choice. The two
automorphisms agree on (E.123.3), since both fix the chosen ninth
root of $-2$. The Galois fibre product therefore supplies their
common extension. Chebotarev, with the finite bad set removed,
gives a new prime $q$ satisfying all these conditions and the
original size bounds. The conjugacy class formulation causes no
problem: fixing some radical root and having the specified local
point-section class are conjugacy-invariant conditions.

The number $b$ is an $n$th power modulo $q$. Choose by Hensel
a unit $z_q$ with $b z_q^n=1-aq^n$, as in E.102.
The cyclotomic congruence gives $v_3(q^2-1)=1$.

\emph{Choosing the new $B$ prime.}
Prescribe $r\equiv\eta_p\pmod{p^{L_p}}$ at every old place
and $r\equiv z_q\pmod q$. The CRT class modulo $Mq$ is
$-1\pmod4$, and its three-adic part has
$v_3(r+1)=2$. Thus it is $-1\pmod{12}$ and fixes $\sqrt3$.
The coefficient $c=aq^n$ has the form $3w^8$ and its square
root is fixed. In $F=\mathbb Q(\mu_{Mq})$ prescribe the
extension fixing a chosen $m$th root of $c$. On $L_9F$ prescribe
\[
 (a_0(F_r),b_0(F_r),h_0(F_r))
       =(0,0,\eta(\chi(F_r)))\pmod9,
 \qquad \eta(u)=\frac{u^2-1}{24}.
 \tag{E.123.8}
\]
The value of $\eta$ modulo nine is specified by the cyclotomic
CRT data: $F$ contains both $\mu_{27}$ and $\mu_8$, so the
integer quotient is determined modulo nine. In particular it is
not being inferred from the multiplier modulo nine alone.
By (E.123.4) the only common radical condition in $L_9F$ is
fixing $3^{1/9}$, which is precisely its first zero abelian
coordinate. The central coordinate in (E.123.8) remains free.
Chebotarev now supplies $r>q$ satisfying all the conditions.
Its Frobenius fixes a root of $X^m-c$, so $c$ is an $m$th
power modulo $r$, and $v_3(r^2-1)=2$.

Set $a_{j+1}=a q^n$ and $b_{j+1}=b r^n$.
The verification of every invariant is unchanged from E.102:
at old places the exact chosen old unit equation makes the new
sum error divisible by $p^{L_p}$; at $q$ the residue $z_q$
gives a simple unit solution; at $r$ the fixed radical root
gives one; at all remaining primes the same finite-field bound
applies. Complete splitting was used in the old verification
only to obtain these residue roots, which have just been
supplied. The powers $n_j$ still tend to zero in every finite
coefficient quotient, and the separate local unit sequences
converge exactly as in E.102. This proves the asserted ordinary
completed pair and its infinite support.

\emph{The canonical class-two reduction.}
Choose literal integral three-adic Kummer cocycles with
\[
 a\equiv a_0,\qquad b\equiv b_0\pmod9,
 \tag{E.123.9}
\]
where $a_0,b_0$ belong to the rational point three; this is
possible because every new exponent is divisible by nine.
By the proved class-two lifting result for ordinary pairs,
choose an integral $h$ with $\delta h=-a\cup b$.
Then $h-h_0\pmod9$ is a cocycle in $\mathbb Z/9(2)$.
It is unramified outside three. At a new support prime $p$,
$p-1$ is a three-adic unit. The literal local primitive of
the zero-valuation coordinate is integral, and the class-two
tame relation gives $h(T_p)=n_j u_p$ or $-n_j v_p$.
These are zero modulo nine. The point-three section is
unramified there. At two the primitive is unique since
$2-1=1$; (E.123.9) makes its difference from the point-three
primitive divisible by nine, and hence makes the central
inertia difference divisible by nine. At every other place
different from three both abelian inertias vanish; the
integral tame relation $(p-1)h(T_p)=0$ forces $h(T_p)=0$.

The exact coefficient sequence and
$H^2(G_{\mathbb Q,\{3,\infty\}},\mathbb Z_3(2))=0$ show that
every such unramified finite cocycle lifts to integral
$H^1$. That group is $\mathbb Z/3$, generated by the actual
$\eta=(\chi^2-1)/24$. Twisting $h$ by the unique multiple
of this integral class and then making a central coboundary
adjustment therefore gives
\[
                   h\equiv h_0\pmod9
 \tag{E.123.10}
\]
as literal cochains, while preserving (E.123.9). This is also
the unique integral class-two class matching the local point
at two: the restriction of this global order-three torsor to
the local weight-two torsor at two is an isomorphism.

\emph{All support matches and the continuous third lift.}
At an $A$-support $q$, write $b=\delta u$. Since $q-1$ is a
three-adic unit, (E.123.7), (E.123.9), and (E.123.10) give
$u\equiv0$, $\alpha\equiv0$, and $\gamma\equiv0\pmod9$
in a compatible Frobenius representative. Thus
\[
 \frac{\gamma-q\alpha u}{q^2-1}\in\mathbb Z_3.
\]
At a $B$-support $r$ write $a=\delta v$. Equations
(E.123.8)--(E.123.10) give $v\equiv0$, $\beta\equiv0$,
and $\gamma\equiv\eta(F_r)\pmod9$. Therefore
\[
 \frac{\gamma+\beta v-\eta(F_r)}{r^2-1}
                 \in\mathbb Z_3.
\]
These are exactly the two actual boundary-point criteria of
B.24. At two, the same primitive comparison with the rational
point three gives the criterion modulo $2^2-1=3$.
Hence every support place different from three is matched.

Finally apply the proof of B.24 in the three-primary component.
Each normalized third residue is a sum of an integral matching
quotient multiplied by $n_j$ and an integral binomial term
whose valuation tends to infinity with $v_3(n_j)$.
Consequently both normalized residue sequences belong to the
strict $c_0$ group; merely checking separate local vanishing
would not suffice here. B.23 and the vanishing of the fixed
unramified weight-three kernel give an actual continuous
global integral class-three lift. Equivalently one can use
E.66 for that kernel. This completes the construction.
\end{proof}

\begin{result}{E.124}{Proved}{Only the even matching layers remain for the modified support pattern}
Let an ordinary pair have the support pattern constructed in E.123.
For $w\geq3$, fix a global integral pro-three class-$(w-1)$ section which agrees
with the actual local point section at every support place
$q\ne3$, and let $b_w$ be the rank of the weight-$w$ lower
central graded piece. A global class-$w$ lift exists by E.66.
For its possible local boundary matches the following statements
hold.

If $w\geq3$ is odd, every such lift is automatically boundary
matched at all those support places. If $w=2k$ is even, put
$s=1+v_3(k)$. The global twisting group is
$(\mathbb Z/3^s)^{b_w}$, and the local twisting groups are
\[
 \begin{array}{c|c|c}
 \text{support place}&\text{local group}&
     \text{image of the global group}\\ \hline
 2\text{ or }q_j&(\mathbb Z/3^s)^{b_w}&
          \text{the full group}\\
 r_j&(\mathbb Z/3^{s+1})^{b_w}&
          3(\mathbb Z/3^{s+1})^{b_w}.
 \end{array}
\]
The actual localization homomorphisms include the indicated
unit factors. Matching at two fixes the unique possible global
twist. It matches everywhere if and only if the discrepancy
tuple at every fresh support agrees with that pinned twist.
Thus an $A$-support still imposes a compatibility condition;
a $B$-support additionally has one local digit outside the
global twisting image in each Hall coordinate.

In particular, the ghosts of E.123 have boundary-matched
class-three sections and global class-four sections. Every
odd matched layer can be crossed automatically once its
predecessor has been matched. This does not solve the
simultaneous even-layer compatibility conditions.
\end{result}
\begin{proof}
For $q\ne3$ and $w\geq2$, the tame Hochschild--Serre
sequence gives
\[
 H^1(G_{\mathbb Q_q},\mathbb Z_3(w))
       =\mathbb Z_3/(q^w-1)\mathbb Z_3.
 \tag{E.124.1}
\]
Indeed the possible inertia term is
$H^0(\widehat{\mathbb Z},\mathbb Z_3(w-1))=0$:
multiplication by the nonzero integer $q^{w-1}-1$ is
injective on the integral Tate module. The unramified first
group is the displayed Frobenius quotient. Wild inertia
does not contribute to these three-primary coefficients.

Every support place different from three is $-1\pmod3$.
For odd $w$, the integer $q^w-1$ is therefore a three-adic
unit, so (E.124.1) vanishes. Since the actual local point
section has a class-$w$ lift, the local difference between
it and any global lift above the matched predecessor lies
in the zero group. This proves automatic matching at each
such place, without a choice of a global twist.

If $w=2k$, the elementary lifting-the-exponent identity is
\[
 v_3(q^w-1)=v_3(q^2-1)+v_3(k).
 \tag{E.124.2}
\]
E.123 gives the first valuation as one at every $q_j$ and
as two at every $r_j$; it is also one at $q=2$. This proves
the asserted local group orders. By E.67 the global group
in a single Hall coordinate has order $3^s$ and is generated
by
\[
 \eta_w(g)=\frac{\chi_3(g)^w-1}{3^s}.
\]
Its local value at Frobenius is
$(q^w-1)/3^s$. At two and the $A$-supports this is a unit
in the local quotient, hence gives an isomorphism. At a
$B$-support it has valuation exactly one, hence generates
the index-three subgroup. Taking the $b_w$ Hall coordinates
gives the table.

More explicitly, if $d_q$ denotes the local difference
from the point section in a fixed normalization, a global
parameter $\lambda\in(\mathbb Z/3^s)^{b_w}$ changes it by
\[
 d_q\longmapsto d_q+
       \lambda\frac{q^w-1}{3^s}
       \pmod{q^w-1}.
 \tag{E.124.3}
\]
The unit factor at two fixes exactly one $\lambda$. The
conditions at all other places are exactly the vanishing
of (E.124.3) for that same value. At a $B$-support even the
existence of a separately matching local twist requires
the discrepancy to lie in the subgroup divisible by three;
after that, its remaining digits must still agree with
the pinned parameter. This is the concrete specialization
of B.26, not a replacement of its simultaneous condition
by separate local solvability.

Apply the odd assertion with $w=3$ to E.123. Its global
class-three lift is boundary matched. E.66 then supplies
a global class-four lift. Repeating the same argument at
any later odd step proves the stated parity reduction.
The preceding odd lift can vary in its global free
$H^1$ group; nothing here asserts that some choice always
makes the following even discrepancy tuple compatible.
\end{proof}

\begin{result}{E.125}{Proved}{A ghost family with equal global and local even-weight parameters}
There are pairwise distinct primes $q_j,r_j\notin\{2,3,5,11\}$
and integers $n_j=(j+11)!$ for which
\[
 A=\frac5{16}\prod_{j\geq0}q_j^{n_j},\qquad
 B=\frac{11}{16}\prod_{j\geq0}r_j^{n_j}
 \tag{E.125.1}
\]
are an ordinary-local completed unit pair with both coordinates
of infinite support. Every support prime $p$ satisfies
\[
        p\equiv-1\pmod3,\qquad v_3(p^2-1)=1.
 \tag{E.125.2}
\]
The pair has a boundary-matched integral pro-three class-two
section, boundary-matched class-three sections, and global
class-four sections. Boundary matching here concerns every
nonzero-valuation place; three is not a support place.

At every subsequent even layer $w=2k$, above a fixed matched
predecessor, the global twisting group and every support-place
twisting group are $(\mathbb Z/3^{1+v_3(k)})^{b_w}$, and every
localization map is an isomorphism. A global matched lift exists
if and only if the local discrepancy vectors, pulled back by
these isomorphisms, are all equal. Failure is witnessed by two
support places, one of which may always be taken to be two.
At every odd layer matching is automatic once the predecessor
is matched. No compatible all-depth family is asserted.
\end{result}
\begin{proof}
The construction parallels E.123, but both base coefficients
are three-adic units. We check all changes, including the
fixed denominator and the old infinity-type support.

Put $a_0=5/16$, $b_0=11/16$, and at stage $j$ use
$n=n_j=(j+11)!$, $m=n_{j+1}$ and
$H_j(p)=2v_p(n_{j+1})+j+3$. The coefficients have the form
\[
 a=5U^8/16,\qquad b=11V^8/16,
 \qquad (UV,2\cdot3\cdot5\cdot11)=1,
\]
with disjoint numerator supports. Retain unit-coordinate local
solubility of $aX^n+bY^n=1$ and
$v_p(a+b-1)\geq H_j(p)$ for $p\leq B(n)$.
The starting point $(1,1)$ verifies these invariants exactly.
Let $S$ contain the primes at most $B(m)$ and all numerator
and denominator supports. Write $L_p=H_{j+1}(p)$ and take
CRT precision $L_p$ except that at two it is $L_2+4$.
The resulting modulus is denoted $M$.

At two choose $\eta_2=1$ and
\[
 \xi_2^n=(1-b)/a,\qquad \xi_2\in1+4\mathbb Z_2.
\]
The right side differs from one with valuation at least
$H_j(2)+4$, because $v_2(a)=-4$. Thus the root exists and
is near one to more than the original required precision.
At primes dividing $m$, other than three, make the near-one
root choices of E.102. At two this is also compatible with
the choice just made. In particular both roots are one
modulo five and eleven, and $\xi_2\equiv1\pmod4$.

At three choose $\xi_3=2$ and let $Y\in1+3\mathbb Z_3$ be
the unique root
\[
 Y^n=(1-a2^n)/b,\qquad \eta_3=-Y.
\]
Writing $f=v_3(n)$, the coefficients $a,b$ are units and
$H_j(3)>f+1$. Hence, exactly as in (E.123.5) but without
the extra valuation of $a$,
\[
 v_3\left((1-a2^n)/b-1\right)=f+1,
 \qquad v_3(Y-1)=1.
 \tag{E.125.3}
\]
Both old roots therefore have distance from minus one of
valuation exactly one. These choices solve the old unit
equation exactly. At all other old places retain any
unit-coordinate solution as before.

We verify the radical-field input for the new coefficients.
For $c=5W^8/16$ or $11W^8/16$ and $h$ divisible by $n_0$,
the field $\mathbb Q(\mu_h)$ contains the square root of
$c$: the quadratic conductors are five and forty-four.
The element $c$ is not an odd prime power in any abelian
extension containing the required roots of unity. Indeed its
valuation one at five or eleven makes the pure radical have
the indicated prime degree, while its real radical field is
not normal. It is not a fourth power either, because that
would put $5^{1/4}$ or $11^{1/4}$ in an abelian extension,
after multiplication by $2/W^2$. The commutator argument
of E.102 now proves, without change,
\[
 [\mathbb Q(\mu_h,c^{1/h}):\mathbb Q(\mu_h)]=h/2,
 \qquad
 \mathbb Q(\mu_h,c^{1/h})^{\rm ab}=\mathbb Q(\mu_h).
 \tag{E.125.4}
\]
Thus the fixed-root extension criterion of E.123 applies.

Use now the fixed class-two point field $L_{9,t}$ for
$t=5/16$. Its full group is $H_9\rtimes(\mathbb Z/9)^\times$.
The Kummer independence proof uses valuation one at five
for $a_0$ and at eleven for $b_0$; the cyclotomic ramification
indices at those primes are prime to three. The same
commutator argument gives maximal abelian field
$\mathbb Q(\mu_9)$. Over every relevant cyclotomic field
$F$, the cyclic radical intersections are consequently
\[
 \begin{split}
 L_{9,t}F\cap F(b^{1/n})&=F(b_0^{1/9}),\\
 L_{9,t}F\cap F((aq^n)^{1/m})&=F(a_0^{1/9}).
 \end{split}
 \tag{E.125.5}
\]
Each line follows because the intersection is cyclic of
degree at most nine and already contains the indicated
degree-nine line; all preceding and current update exponents
are divisible by nine.

Choose $q$ in the CRT class of $\xi_p$. This class is
$1\pmod{44}$, so fixes $\sqrt{11}$ and the square root
of $b$. Prescribe the fixed-root automorphism on the
$n$th radical field and the point-$t$ coordinates
$(a_t,b_t,h_t)=(0,0,0)\pmod9$. They agree on the first
intersection in (E.125.5), and Chebotarev gives arbitrarily
large new primes $q$. The residue field has an $n$th root
of $b$; choose the same local $z_q$ as in E.123.

Choose $r$ in the old CRT classes $\eta_p$, and in the
class of $z_q$ modulo $q$. Its class is $1\pmod5$, so
fixes $\sqrt5$ and the square root of $aq^n$. Prescribe
the fixed-root automorphism on the $m$th radical field and
the point-$t$ coordinates
\[
              (a_t,b_t,h_t)=(0,0,\eta(\chi))\pmod9.
\]
The cyclotomic modulus includes $\mu_{27}$, so the last
coordinate is well defined. Compatibility is exactly the
second line of (E.125.5). Chebotarev gives a new $r>q$.
Equation (E.125.3) ensures that both new primes satisfy
(E.125.2).

Set $a_{j+1}=a q^n$, $b_{j+1}=b r^n$.
At an old odd prime the original estimates apply literally.
At two, the extra four CRT digits make
$v_2(aq^n+br^n-1)\geq L_2$. For local solubility multiply
the new polynomial by sixteen: its two coefficients are
odd units, its sum error has valuation at least $L_2+4$,
and the derivative in either variable at one has valuation
$v_2(m)$. The same Hensel estimate supplies a unit root.
At $q$ and $r$ the fixed-root conditions give the same
simple residue solutions as before. Every remaining prime
is outside the fixed denominator and the numerator support,
so the finite-field point bound of E.102 applies. This
proves every inductive invariant for rational coefficients.

The completed products and the ordinary local unit parts
converge by the factorial argument. At two the two valuations
are both $-4$; at five they are $(1,0)$, at eleven $(0,1)$,
and at each fresh support they are the specified positive
factorial and zero. At all other primes they are zero.
The sum-error invariant proves $A_p+B_p=1$ for their
ordinary values at every finite prime. Both signs are
positive at the real place and are represented by the
ordinary real pair $(1/2,1/2)$. The infinite disjoint
supports prove that the completed coordinates are not
rational Kummer classes.

To normalize a global integral class-two lift modulo nine,
use literal $a\equiv a_t$, $b\equiv b_t\pmod9$.
At every fresh support, and at the old places five and
eleven, the primitive division by $p-1$ is by a three-adic
unit. The same inertia argument as in E.123 applies.
For completeness, the exact general class-two tame relation
in these coordinates is
\[
 (p-1)h(T)=m\beta-n\alpha+(p-1)mn/2,
 \qquad (m,n)=(a(T),b(T)).
 \tag{E.125.6}
\]
It follows by comparing the central coordinates in
$F T F^{-1}=T^p$. At two, $m=n=-4$ and hence
$h(T)=4(\alpha-\beta)+8$. Its difference from the
point-$t$ value is divisible by nine. Thus $h-h_t\pmod9$
is unramified outside three at the old infinity-type
support as well. E.66 and the integral weight-two
order-three class normalize it to zero, exactly as before.

The new support matches now follow from their prescribed
Frobenius coordinates and denominators of valuation one.
At two, five, and eleven, every local unit has trivial
three-primary Kummer class, since its residue unit group
has order prime to three. The ordinary local pairs have
the same abelian classes as the base point. The proved
local boundary singleton statement A.31 therefore
identifies their actual point-section classes with that
of the base point. To align literal abelian cocycles,
the conjugating abelian coordinates are divisible by
nine, because their Frobenius divisions are by $p-1$.
The class-two Hall formulas have only the denominator
two, so their central changes are divisible by nine
as well. The global central congruence therefore makes
the local difference zero in the order-three central
torsor. All old supports are matched.

The strict convergence proof of B.24 applies to the
fresh factorial supports; the finitely many old matched
places add no convergence issue. It gives a global
class-three lift. Odd local first cohomology vanishes
at every support by (E.125.2), so every such lift is
matched, and E.66 gives a global class-four lift.

Finally (E.124.1)--(E.124.3) apply with
$v_3(p^2-1)=1$ at every support. Thus the local map
from the global even-weight torsor is an isomorphism
at every one of them. Pulling the discrepancy vectors
back by these maps gives a family $\lambda_p$ in one
fixed finite group. It is the image of a single global
twist exactly when $\lambda_p$ is independent of $p$.
If not, some $\lambda_p\ne\lambda_2$ supplies the
claimed two-place witness. This is a sharper form of
the remaining gate, not a proof that the gate vanishes.
\end{proof}

\begin{result}{E.126}{Proved}{Full geometric image with the actual higher Galois action}
Let $\Pi$ be the free geometric pro-three fundamental
group of $\mathbb P^1-\{0,1,\infty\}$ at the integral
tangential base point $\overrightarrow{01}$. Let
\[
 \rho:G_{\mathbb Q}\longrightarrow\operatorname{Aut}(\Pi)
\]
be its actual Galois action, put $\mathcal A=\rho(G_{\mathbb Q})$,
and let $z_t:G_{\mathbb Q}\to\Pi$ be the point cocycle of
$t=5/16$, after choosing a geometric path. Then
\[
 z_t(\ker\rho)=\Pi,
 \qquad
 \{(z_t(g),\rho(g)):g\in G_{\mathbb Q}\}
                  =\Pi\rtimes\mathcal A.
 \tag{E.126.1}
\]
In particular every finite quotient retaining the actual
arithmetic action has the full geometric image. Prescribing
finite geometric Frobenius coordinates is therefore possible
above every arithmetic element which occurs, before any
additional radical-field intersection conditions are imposed.
This statement does not replace the higher arithmetic action
by its cyclotomic character.

The same finite-quotient conclusion holds for any global
integral class-$r$ section whose abelian valuation vectors
at five and eleven are $(1,0)$ and $(0,1)$, respectively.
In particular it applies to every such section above the
rational pair $(5/16,11/16)$ or an E.125 ghost. It also
holds after adjoining any finite arithmetic data
unramified outside three, including the cusp path values.
\end{result}
\begin{proof}
The actual based action $\rho$, not merely its outer class,
is unramified outside three. This is the direct specialization
of \citet[Definition 2.0 and Theorem 2.1, pp.198--200]
{AuditWojtkowiak2007}. That definition explicitly allows
tangential points. For the punctures $0,1,\infty$ and the
unit tangent from zero towards one, every finite prime
has good reduction in the stated definition; the only
excluded prime for the pro-three action is consequently
three. Thus the inertia groups at five and eleven lie
in $\ker\rho$.

On $\ker\rho$ the cocycle identity for $z_t$ is an
ordinary homomorphism identity. Its image $K$ is a closed
subgroup of $\Pi$, because the kernel is compact and
$z_t$ is continuous. The two abelian coordinates of the
point cocycle are the Kummer characters of $5/16$ and
$11/16$. A three-primary tame inertia generator at five
therefore has image $(1,0)$ in
$\Pi/\Phi(\Pi)\simeq\mathbb F_3^2$, and one at eleven
has image $(0,1)$. These assertions require only their
ordinary valuations; the usual local peripheral
description also identifies the full inertia images
with conjugates of the two cusp generators.

It follows that $K\Phi(\Pi)=\Pi$. For clarity, the
Frattini implication used here needs no density
assumption: if a closed subgroup of a finitely generated
pro-three group were proper, its image would be proper
in some finite three-group quotient. A maximal subgroup
of that quotient has index three, and its preimage is
the kernel of a nonzero homomorphism to $\mathbb F_3$.
Such a kernel contains the Frattini subgroup and cannot
contain a subgroup surjecting onto the Frattini quotient.
This contradiction proves $K=\Pi$.

The joint image is a subgroup of $\Pi\rtimes\mathcal A$,
contains $\Pi\times\{1\}$ by the first equality, and
projects onto $\mathcal A$ by its definition. Multiplying
any element above a specified arithmetic image by an
element of the full geometric kernel gives every element
of that fibre. This proves the second equality, and
passing to finite quotients proves the last assertion.

For the stated generalization, replace the point cocycle
by the given section cocycle in its finite geometric
quotient, and replace the kernel of $\rho$ by the kernel
of all the chosen arithmetic data. Inertia at five and
eleven still lies in this kernel. Its two images on the
Frattini quotient are still the same Kummer valuation
vectors. The compact-image and Frattini argument is
unchanged. Thus no point property of the section is
needed for this finite-image conclusion.

The result supplies finite Frobenius freedom. It does
not assert a uniform common prime satisfying infinitely
many Frobenius conditions, nor a compatible system of
boundary-matched global ghost sections.
\end{proof}

\begin{result}{E.128}{Proved}{Finite Frobenius prescriptions and fixed radical roots are compatible}
Let $L/\mathbb Q$ be a finite Galois extension containing
$\mu_3$ such that $L/\mathbb Q(\mu_3)$ is a three-extension
and $L^{\rm ab}$ is contained in a cyclotomic field of
three-power conductor. Let $h$ and $c$ be radical data from
E.123 or E.125, so $24\mid h$,
\[
 [\mathbb Q(\mu_h,c^{1/h}):\mathbb Q(\mu_h)]=h/2,
 \qquad
 \mathbb Q(\mu_h,c^{1/h})^{\rm ab}=\mathbb Q(\mu_h).
\]
Let $F$ be a cyclotomic field containing $\mu_h$ and
$L^{\rm ab}$. Suppose $\sigma\in\operatorname{Gal}(F/\mathbb Q)$
and $g\in\operatorname{Gal}(L/\mathbb Q)$ agree on $L\cap F$,
and
\[
 \chi_3(\sigma)\not\equiv1\pmod3,
 \qquad \sigma(\sqrt c)=\sqrt c.
 \tag{E.128.1}
\]
Then there are arbitrarily large primes in the prescribed
cyclotomic CRT class whose Frobenius on $L$ is conjugate to
$g$ and for which $c$ is an $h$th power in the residue field.
Any prescribed finite set of primes can be omitted.

In the family of E.125 this applies to arbitrary finite
point-section Frobenius conditions supplied by E.126,
including the actual finite cusp-one path cocycle as part
of the arithmetic data. Thus any fixed finite cusp-zero
or cusp-one matching condition can be imposed at the
corresponding new prime without losing the radical-root
or CRT requirements. This is a finite extension lemma;
it asserts no prime satisfying infinitely many conditions.
\end{result}
\begin{proof}
The two given automorphisms extend to one automorphism
$\widetilde g$ of $LF$, by their agreement on $L\cap F$.
Put $R=F(c^{1/h})$. The radical maximal-abelian calculation
shows that $R/F$ is cyclic of degree $h/2$.
The intersection $I=LF\cap R$ is therefore cyclic over $F$.
Since $F$ contains $\mu_3$ and $L/\mathbb Q(\mu_3)$ is
a three-extension, $LF/F$ is a three-extension. Consequently
\[
             I=F(c^{1/3^d})
 \tag{E.128.2}
\]
for some $0\leq d\leq v_3(h)$. Here the cyclic radical
extension has a unique subfield of each three-power degree;
the indicated subfield has degree $3^d$ because $c$ is not
a cube in the cyclotomic field $F$.

Choose any $\delta$ with $\delta^{3^d}=c$ and a primitive
$3^d$th root of unity $\zeta$. On these roots the action of
$\widetilde g$ has the form
\[
 \delta\zeta^k\longmapsto\delta\zeta^{a+u k},
 \qquad u=\chi_3(\sigma)\pmod{3^d}.
\]
The congruence $a+(u-1)k=0$ has a unique solution, since
$u-1$ is a three-adic unit. Thus $\widetilde g$ fixes one
of the roots $\delta$. This is the step for which the
minus-one residue class matters; complete radical splitting
would require a different compatibility argument.

Choose an $h$th root $z$ of $c$ with
\[
 z^{h/3^d}=\delta,\qquad z^{h/2}=\sqrt c,
 \tag{E.128.3}
\]
where the square root is the one fixed by $\sigma$.
These choices are compatible: multiplication of an arbitrary
$h$th root by $\zeta_h^k$ independently prescribes its
$3^d$-root and square-root images, by the Chinese remainder
theorem for $3^d$ and two. As in E.123, the irreducible
polynomial $X^{h/2}-\sqrt c$ over $F$ shows that $\sigma$
extends to an automorphism of $R$ fixing $z$. It agrees with
$\widetilde g$ on $I$, because they agree on $F$ and both
fix $\delta$. The Galois fibre product therefore supplies
an automorphism of $LFR$ with all the desired restrictions.
Chebotarev supplies its conjugacy class at arbitrarily
large unramified primes. A conjugate still fixes some
$h$th root of $c$, so the residue polynomial has a root;
the cyclotomic restriction is unchanged by conjugation.
This proves the first assertion, without requiring a
calculation of the exact exponent $d$.

We justify the asserted supply of fields and automorphisms
for E.125. Fix a finite geometric three-group quotient $P$
and a finite arithmetic data field $D$, unramified outside
three, retaining the actual action on $P$ and every finite
cusp path value used in the prescription. Such a field
exists. The actual based action is unramified outside
three by E.126; the tangential cusp path torsor has the
same good-reduction property by
\citet[Definition 2.2 and Proposition 2.3, p.200]
{AuditWojtkowiak2007}. If necessary, retain the cusp-one
path modulo one higher class than the action: its last
central coordinate can disappear from conjugation in
the preceding class and must not be discarded.

Adjoin $\mu_3$ to $D$. Its Galois group over
$\mathbb Q(\mu_3)$ is a three-group. Indeed the arithmetic
action on the geometric Frattini quotient is the scalar
cyclotomic character; the kernel of the automorphism
action on that quotient is pro-three, and adding finite
pro-three path coordinates adds only a three-group kernel.
One can verify the automorphism assertion by the usual
lower exponent-three filtration: an automorphism trivial
on the Frattini quotient has successive displacement
quotients contained in homomorphism groups between
finite elementary abelian three-groups, so every finite
image of that kernel has three-power order.

Let $L$ retain this arithmetic data and the point-$t$
cocycle in $P$. The proof of E.126 works with $\ker(G_{\mathbb Q}
\to\operatorname{Gal}(D/\mathbb Q))$ in place of $\ker\rho$:
inertia at five and eleven is still in that kernel and
still gives the geometric Frattini basis. Hence
\[
 \operatorname{Gal}(L/\mathbb Q)
      =P\rtimes\operatorname{Gal}(D/\mathbb Q).
\]
In the abelianization, complex conjugation acts on
$P^{\rm ab}$ by minus one. Since $P^{\rm ab}$ is a
three-group, all its coinvariants are zero. Thus $P$
lies in the commutator subgroup, and $L^{\rm ab}$ is
an abelian subfield of $D$. Kronecker--Weber and its
unramifiedness outside three put it in a cyclotomic
field of three-power conductor. The hypotheses of
the first assertion are satisfied.

Finally, choose an actual arithmetic element with the
desired three-adic cyclotomic value, namely the selected
old local root at three. The cyclotomic character is
surjective, and all abelian restrictions in $D$ are
determined by that value. The full geometric image then
allows its geometric coordinate to be the finite cusp-zero
coordinate, or the actual finite cusp-one path coordinate.
The resulting $g$ agrees with the CRT automorphism on
$L^{\rm ab}$. Enlarge only the finite three-power CRT
precision to contain that field, keeping the chosen
old roots. The existing quadratic congruences in E.125
give the second condition of (E.128.1), while its prime
shells give the first. The first part of the proof now
imposes the finite marked condition simultaneously with
the original radical root and every old-place congruence.

By the final assertion of E.126, the same proof works
when the fixed finite field is obtained from any global
integral class-$r$ section with the two prescribed
valuation vectors at five and eleven. In particular,
finite class-four representatives above the rational
base pair need not be the actual point representative.
Their arithmetic data field is still unramified outside
three; the full geometric image supplies the same
finite cusp prescriptions and the identical radical
compatibility argument applies.
\end{proof}

\begin{result}{E.129}{Proved}{Selected boundary-matched ghosts still fail the compact pro-three gate}
There are ordinary infinite-support pairs with every property
of E.123--E.124, including boundary-matched integral pro-three
class-two and class-three sections and global class-four
sections, which have no full pro-three section.
They are obtained by varying the first old root at three
within the same valuation shells and prescribing one extra
finite Frobenius class at the first $B$-support prime.
The extra condition makes
\[
 P_0-T=(3,-6)\in E'(\mathbb F_{r_0}),
 \qquad E':Y^2=X^3+9,
\]
have three-primary order three, while
$E'(\mathbb F_{r_0})\{3\}$ has order nine.

An exact auxiliary witness for the simultaneous finite
conditions is the prime $6551$. This is a witness for a
Frobenius class, not a proposed actual recursive support
prime. The constructed $r_0$ satisfies all the much larger
CRT and radical requirements and may be arbitrarily large.
The result excludes selected pairs, not the whole family
of E.123 and not the family of E.125.
\end{result}
\begin{proof}
We first identify the actual marked ninth-Heisenberg test,
then give its finite witness and extend that witness to
the ordinary recursion.

\emph{The marked geometric quotient and its boundary value.}
Over $K=\mathbb Q(\zeta_9)$ put
\[
 \Theta(X)=\prod_{i=1}^{8}(1-\zeta_9^iX)^i,
 \qquad
 X^9=t,\quad Y^9=1-t,\quad Z^9=\Theta(X).
 \tag{E.129.1}
\]
Over $K(X)$ the two Kummer classes $1-X^9$ and
$\Theta(X)$ are independent modulo ninth powers:
the valuations at $X=1$ and $X=\zeta_9^{-1}$ of a
product $(1-X^9)^a\Theta(X)^b$ force first
$a\equiv0$ and then $b\equiv0\pmod9$.
Thus the degree over $K(t)$ is $9^3$.

The two deck lifts are
\[
 x:(X,Y,Z)\longmapsto
     (\zeta_9X,Y,Z(1-X)/Y),
 \qquad
 y:(X,Y,Z)\longmapsto(X,\zeta_9Y,Z).
 \tag{E.129.2}
\]
The identity
$\Theta(\zeta_9X)/\Theta(X)=((1-X)/Y)^9$
checks the first lift. Both have order nine, and their
commutator acts on $Z$ by multiplication by $\zeta_9$.
They generate the degree-$9^3$ deck group $H_9$.
These are the actual peripheral markings: near the
tangential base point at zero, take the branches
$Y=1+O(X^9)$ and $Z=1+O(X)$; the first lift is the
monodromy of the $X$ root. Near the branch $X=1$,
$\Theta(X)$ is a unit and the second lift is the
monodromy of the $Y$ root with $Z$ unchanged. Hence
(E.129.1) is the actual class-two exponent-nine
quotient with its zero and one markings, rather than
an unmarked cover of the same degree.

Suppose $p\equiv-1\pmod9$, and work over
$\mathbb F_{p^2}$ with a primitive $\zeta_9$.
For the based arithmetic action with cyclotomic
multiplier $-1\pmod9$, conjugation of the coefficients
gives
\[
 \Theta(X)^{\sigma}=
       \frac{H(X)^9}{\Theta(X)},
 \qquad H(X)=\frac{1-X^9}{1-X}.
\]
The normalized branch $Z(0)=1$ therefore transforms
as $Z\mapsto H(X)/Z$. This fixes the based arithmetic
normalization, not only the geometric deck action.

At $t=3$, the ninth roots $X$ of $3$ and $Y$ of
$-2$ lie uniquely in $\mathbb F_p$, since
$(9,p-1)=1$. The remaining central Frobenius factor
is
\[
 C_p(3)=\frac{\Theta(X)^{(p+1)/9}}{H(X)}\in\mu_9.
\]
At the tangential branch over one, $X=1$, $Y=0$,
and $H(1)=9$. Its actual central Frobenius factor is
\[
 C_p(1)=\frac{\Theta(1)^{(p+1)/9}}9.
\]
Consequently the exact equality of the point-three
and cusp-one class-two Frobenius classes is
\[
 \mathcal R_p=
    \left(\frac{\Theta(X)}{\Theta(1)}\right)^{(p+1)/9}
          \frac9{H(X)}=1.
 \tag{E.129.3}
\]
The factors are independent of the chosen ninth roots
of $Z^9$, because $9\mid p+1$.

The cusp term also checks the Ihara normalization
explicitly. Pairing $i$ with $9-i$ gives
$\Theta(1)=\pm3^9\zeta_9^3$. For example, if
$U=\prod_{i=1}^4(1-\zeta_9^i)$, then
$U^2=9\zeta_9$ and
$\Theta(1)=\zeta_9^{-60}U^9$.
Since $(p+1)/9$ is even,
\[
 C_p(1)=\zeta_9^{(p+1)/3}
       =\zeta_9^{-\eta(F_p)},
 \qquad \eta(F_p)=\frac{p^2-1}{24}\pmod9.
 \tag{E.129.4}
\]
The last equality uses $p\equiv-1\pmod9$ and the
integral three-adic interpretation of the quotient.
Thus the cusp factor is exactly the actual shifted
boundary class used in B.24 and E.123. The negative
exponent is the contravariant function-field coordinate
for the central Hall exponent. In particular, after
aligning the zero abelian Frobenius coordinates,
(E.129.3) is precisely $\gamma\equiv\eta(F_p)\pmod9$.

\emph{An exact finite witness.}
Take $p=6551\equiv71\pmod{108}$ and
$\mathbb F_{p^2}=\mathbb F_p[s]/(s^2-17)$.
The following values, each reduced modulo $p$, give
an independently reproducible certificate:
\[
 \begin{array}{c|c}
 \zeta_9&6318+4878s\\
 X\ (X^9=3)&5423\\
 H(X)&940\\
 \Theta(X)&2167+4130s\\
 \Theta(1)&6536+617s\\
 C_p(3)=C_p(1)&3275+1508s=\zeta_9^6.
 \end{array}
 \tag{E.129.5}
\]
Here $17$ is a quadratic nonresidue, $\zeta_9$ has
exact order nine, and $\eta(F_p)=3\pmod9$.
The last equality verifies (E.129.3), including the
nonzero cusp correction. On $E'$ the same finite
calculation gives
\[
 \left[\frac{6551+1}{9}\right](3,-6)
       =[728](3,-6)=(0,-3)\ne O.
 \tag{E.129.6}
\]
This point has order three.
For every $p\equiv-1\pmod3$, cubing is a bijection
of $\mathbb F_p$, so counting the equation
$Y^2=X^3+9$ gives $\#E'(\mathbb F_p)=p+1$.
Its rational three-torsion consists of $O,(0,3),(0,-3)$:
the other division-polynomial possibility is
$X^3=-36$, with $Y^2=-27$, and $-3$ is not a square
when $p\equiv-1\pmod3$. Thus its three-primary
subgroup is cyclic. At the present prime its order
is nine. Equation (E.129.6) says that $P_0-T$ belongs
to $3E'(\mathbb F_p)$ but not to $9E'(\mathbb F_p)$.

The actual code is
\path{code/ghost_audit_compact_h9.py}.
Its recorded output is
\path{code/ghost_audit_compact_h9.out}.
Besides this certificate, the code checks both
cyclotomic norm identities used in (E.129.3)--(E.129.4).
Its auxiliary search is not used in the infinite
existence argument.

\emph{The finite joint field.}
Let $M$ be the Galois field of all points
$[9]^{-1}(P_0-T)$ on $E'$, retaining $E'[9]$,
and let $L$ be its compositum with the point-three
field $L_9$ and $\mathbb Q(\mu_{27})$.
It is unramified outside two and three. For $M$
this follows from the good elliptic model and the
finite etale multiplication-by-nine torsor away
from those primes; the point extends over every
such local ring by properness. For $L_9$ it follows
from the marked point and tangential action.

The extension $L/\mathbb Q(\mu_3)$ is a
three-extension. For the elliptic part, the rational
three-torsion point and the determinant one over
$\mathbb Q(\mu_3)$ put the mod-three torsion image
in the order-three upper unipotent subgroup.
The kernel of reduction from mod-nine matrices is
a three-group, and the division torsor adds only
translations by $E'[9]$. The Heisenberg and
cyclotomic parts are three-extensions as well.
It follows that $L^{\rm ab}$ has its only possible
quadratic subfield in $\mathbb Q(\mu_3)$.
Its three-primary abelian subextension is unramified
at two: neither wild two-inertia nor the tame
abelian quotient $\mathbb F_2^\times$ contributes
three-torsion. Kronecker--Weber therefore puts
$L^{\rm ab}$ in a cyclotomic field of three-power
conductor.

Let $g$ be a Frobenius at $6551$ in $L$.
The certificate gives its boundary-matching
Heisenberg class. Its action on the finite torsor
$[9]^{-1}(P_0-T)$ has no fixed point, by (E.129.6).
These are conjugacy-invariant assertions.

\emph{Extending the witness through the recursion.}
At the first stage of E.123, keep $n=7!$ but choose
the old three-adic root $\eta_3=6551$. Its distance
from minus one has valuation two. Let $X_3$ be the
unique root in $1+3\mathbb Z_3$ of
\[
 X_3^n=\frac{1+2\eta_3^n}{3},
 \qquad \xi_3=-X_3.
\]
If $f=v_3(n)$, the expression on the right differs
from one with valuation $f+1$, because
$v_3(\eta_3^n-1)=f+2$. Hence $v_3(X_3-1)=1$.
These roots satisfy $3\xi_3^n-2\eta_3^n=1$
exactly and retain the two shells of E.123.
All other old roots are chosen as in that proof.

Enlarge the finite three-power CRT precision so
that its cyclotomic field contains $L^{\rm ab}$.
Choose $q_0$ by the same fixed-root and cusp-zero
Heisenberg conditions as in E.123. For $r_0$ use
the old root $\eta_3=6551$, the other old roots,
and the required $z_{q_0}$ residue. The resulting
cyclotomic automorphism agrees with $g$ on
$L^{\rm ab}$: its three-adic part is the actual
Frobenius value $6551$. It fixes $\sqrt{3q_0^n}$
by the unchanged residues $-1$ modulo four and
three. E.128, applied to $c=3q_0^n$ and
$h=n_1=8!$, therefore supplies arbitrarily large
$r_0>q_0$ in this CRT class, fixing a radical root
and having the prescribed Frobenius on $L$.

The latter condition imposes the same actual
cusp-one Heisenberg match as E.123, but also makes
$P_0-T\notin9E'(\mathbb F_{r_0})$.
It keeps $v_3(r_0+1)=2$ by its cyclotomic
restriction modulo twenty-seven. Thus
$E'(\mathbb F_{r_0})\{3\}$ is cyclic of order
nine and $P_0-T$ has nonzero three-primary part.
Its divisibility by three follows either from the
same Frobenius on the third-division sub-torsor or
from the Heisenberg boundary-fibre condition.
Continue all subsequent stages exactly as in E.123.
Only an admissible first pair of old roots was
changed, so every ordinary convergence and support
matching proof remains valid. E.124 supplies the
matched class-three and global class-four lifts.

\emph{Excluding a full pro-three section.}
If a full pro-three section existed, its class-two
reduction modulo three would be the fixed point-three
one, by the unique local pin at two and the
normalization in E.123. It would therefore factor
through the actual pointed $H_3$ cover of E.110.
Its proper component would be
$aP_0+bT\in E'(\mathbb Q)^{\wedge3}$ by E.112.
At $r_0$ the actual local point lies over the
cusp-one tube. The boundary identity in E.113 is
$g_0=[3]+P_0$, and that boundary has image $T$.
Reduction and the proved local boundary rigidity
therefore give
\[
                   (3a+1)P_0=T
          \quad\hbox{in }E'(\mathbb F_{r_0})^{\wedge3}.
\]
Multiplying by three shows $3P_0=0$, since
$3a+1$ is a three-adic unit. The same equation
then says $P_0=T$. This contradicts the nonzero
three-primary part of $P_0-T$ imposed above.
All ingredients of the compact obstruction apply
to this modified recursion: the fixed marked
cover and the canonical mod-three class were
preserved. No assertion about all other outputs
has been made.
\end{proof}

\begin{result}{E.133}{Proved}{A dense fourth-layer matching exclusion after every finite prefix}
Every legal finite prefix of the E.125 recursion has an
infinite legal continuation whose distinguished class-two
section admits no global integral class-four lift matched
to the ordinary local point sections at every support
prime. In fact, after imposing the match at two, all nine
possible odd-layer parameter classes are excluded at one
new $A$-support prime.

Every such output still has boundary-matched class-two
and class-three sections and global class-four sections,
as proved in E.125. It has no full pro-three section.
The failure occurs at two support places different from
the coefficient prime three; no local point-section
condition at three is used.

In the topology of legal finite prefixes, the locus with
no support-matched fourth-layer lift is open and dense.
Its complement is a nonempty closed nowhere dense set,
by E.132. More precisely, the consistent prefixes of
E.132 have both a matched fourth-layer continuation and
a continuation with the failure described here.
This describes one matching layer, not an all-depth
survival theorem.
\end{result}
\begin{proof}
We use the characteristic quotient $Q$ and the uniform
finite normalization of E.130--E.131. In particular
\[
 V=\Gamma_4Q\simeq\mathbb F_3^3,
 \qquad \lambda\in\mathbb F_3^2,
 \qquad \tau=-B_2\lambda
 \tag{E.133.1}
\]
are the top central group, the nine odd-parameter
classes, and the unique even twist prescribed at two.
Choose the nine rational-base integral class-four
sections
\[
 s_\lambda=s_0+j_{0,1}\lambda_1+j_{0,2}\lambda_2
                   +\eta_4(-B_2\lambda)
\]
as in (E.132.1), using fixed integral representatives
of the finite parameters. Their reductions in $Q$
include every possible finite reduction of a global
class-four lift above an E.125 pair which matches at
two. This conclusion is the integral normalization
of E.130, not an assumption that arbitrary finite
cocycles lift.

\emph{A common arbitrary central displacement.}
Choose a finite arithmetic data field $D$, unramified
outside three, containing $\mu_9$, the actual action
on $Q$, the finite cusp path values used in E.131,
and the finite odd cocycle $\gamma\bmod3$.
Write $G_D=\operatorname{Gal}(\overline{\mathbb Q}/D)$.
On this group the reference point cocycle $z_0$
is a homomorphism to $Q$. It is surjective by the
finite-class version of E.126: inertia at five and
eleven is in $G_D$ and gives the two geometric
Frattini basis vectors.

On $G_D$ each difference from $s_0$ to $s_\lambda$
has zero odd projection, since $\gamma\bmod3=0$.
It therefore takes values in $V$. The action on $V$
is trivial there: $D$ contains $\mu_3$ and $V$ has
weight four. Each such difference is consequently
an ordinary homomorphism $G_D\to V$.
In particular it vanishes on
\[
                         C=\overline{[G_D,G_D]}.
\]
Surjectivity of $z_0$ gives $z_0(C)=Q'$.
Indeed commutators of arbitrary lifts of elements
of $Q$ map to their commutators, and those generate
the finite commutator subgroup. Since $V\subset Q'$,
every $u\in V$ has a lift $h\in C$ such that
\[
                z_0(h)=z_\lambda(h)=u
                         \quad\hbox{for all }\lambda.
 \tag{E.133.2}
\]
The same argument applied directly to each
$j_{0,i}\bmod3$ says $j_{0,i}(h)=0$.
The element $h$ fixes every cyclotomic field,
because $C\subset\overline{[G_{\mathbb Q},G_{\mathbb Q}]}$.
Thus this common displacement changes neither
the arithmetic data nor any CRT condition.

Let $g_0$ be a prospective Frobenius element in the
new $A$-prime CRT class, with its reference section
equal to the actual finite cusp-zero section in
$Q$. Such an element exists by E.126 and E.128,
with any required finite cyclotomic precision.
More explicitly, first choose it in the compositum
of the reference-point, cusp, and arithmetic data
field with that cyclotomic field; their intersection
is the already prescribed abelian one. Then extend
this element to the full joint finite field retaining
the $j_{0,i}$. No values of the latter cocycles have
been prescribed at this stage, so this extension
adds no compatibility condition.
For the moment retain the restrictions of an
actual lift of $g_0$ in the absolute Galois group;
the values of the additional $j_{0,i}$ are then
defined as well. The lower $H_9$ condition is the
one prescribed in E.125.

For a fresh support, E.131 expresses the residual
matching equation after the pin at two as
\[
 \mathcal E_\lambda(g)=P_4(g)+Z(g)\lambda
                           -k(g)B_2\lambda
                         \quad\text{in }V,
 \qquad k(g)=\frac{\chi(g)^4-1}{3}\pmod3.
 \tag{E.133.3}
\]
Here $k(g)$ is a unit and depends only on the
fixed cyclotomic data. The target cusp and all
alignments are those in E.131, including its
finite-to-integral comparison.
Replace $g_0$ by $g=g_0h$, with $h$ as in
(E.133.2). The cusp and arithmetic data are
unchanged. The actual action of $g_0$ on $V$
is multiplication by $\chi(g_0)^4=1\pmod3$.
Consequently
\[
 P_4(g)=P_4(g_0)+u,
 \qquad Z(g)=Z(g_0),\qquad k(g)=k(g_0),
\]
and therefore
\[
              \mathcal E_\lambda(g)
                   =\mathcal E_\lambda(g_0)+u
                   \quad\hbox{for all nine }\lambda.
 \tag{E.133.4}
\]
There are at most nine forbidden values
$-\mathcal E_\lambda(g_0)$ in the group $V$ of
order twenty-seven. Choose any of the remaining
at least eighteen values for $u$. Then no
$\mathcal E_\lambda(g)$ is zero.
The displacement lies in $\Gamma_4Q$, so the
prescribed $H_9$ Frobenius condition is unchanged.

\emph{Realizing the displacement in the legal recursion.}
Let $L$ retain the finite arithmetic data, all nine
section representations, and the finite cocycles
used above. It is a finite Galois extension,
unramified outside $\{2,3,5,11\}$, by E.131.
It is a three-extension over $\mathbb Q(\mu_3)$:
all retained geometric kernels are three-groups,
and the actual arithmetic data has the same
property as in E.128.
Its only possible quadratic subfield is therefore
$\mathbb Q(\mu_3)$. Its three-primary abelian
subextension is unramified at two, five, and
eleven, since wild inertia has the wrong prime
and the tame abelian inertia orders divide
$1,4,10$, respectively. Kronecker--Weber puts
$L^{\rm ab}$ in a three-power cyclotomic field.

Keep every earlier prime and local choice in the
given finite prefix. At the next stage choose
the old local roots as in E.125 and enlarge
only finite three-power CRT precision to include
$L^{\rm ab}$. The element $g_0$ above can be
chosen in this CRT class. The common displacement
$h$ fixes its entire cyclotomic field, so $g$
has the same restriction. The square-root
condition for the new $A$ radical is still
the unchanged residue $1\pmod{44}$; its
three-adic multiplier is still minus one modulo
three. E.128, applied to $L$ and the current
radical $b_j^{1/n_j}$, gives an arbitrarily
large new prime $q_j$ with this Frobenius class
and an $n_j$th residue root of $b_j$.

Choose the partner $r_j$ by the usual legal
E.125 construction and continue indefinitely.
The central displacement has not changed any
class-two Frobenius condition. Hence all the
ordinary limits, strong class-two matches,
automatic class-three matches, and global
class-four existence assertions of E.125
continue to hold. At the chosen new prime,
however, (E.133.4) excludes every class-four
candidate which matches at two. By the uniform
finite representation in E.131 these equations
are unchanged by every later tail.

Both witnessing places, two and $q_j$, are
different from three. A full pro-three section
would match the actual local points there by
A.31, since both places have a nonzero ordinary
valuation. Its class-four image would contradict
the preceding exclusion. No assertion that the
section is a point section at the coefficient
prime is involved.

Finally E.69 and E.131 make inconsistency an
open finite-prefix condition. The construction
just given makes that condition hold after
every legal prefix, so the excluded locus is
dense. E.132 gives a nonempty closed complement
and a matched continuation for every consistent
prefix. Density makes that complement nowhere
dense. These topological conclusions retain
their finite-layer quantifiers throughout.
\end{proof}

\begin{result}{E.134}{Proved}{The compact gate of E.123 is already necessary at the fourth matching layer}
For an E.123 pair, the existence of a global integral
pro-three class-four section matched to the ordinary
local point sections at every support prime implies
all compact support tests of E.113. More precisely,
the $A$-support tests already follow from its class-two
matching. If a $B$-support prime $r$ fails its compact
test, no global integral class-four lift of the
distinguished class-two section can match at both
two and $r$.

The detecting geometric quotient is the actual marked
triangle quotient of order $243$ and nilpotent class
four. Its kernel over the pointed $H_3$ quotient is
the Galois module $E'[3]$. Consequently the selected
outputs of E.129 fail the fourth matching layer,
although they possess global class-four sections.
No local point-section condition at three is imposed.
Passing these compact tests is not asserted to be
sufficient for fourth-layer matching.
\end{result}
\begin{proof}
\emph{The marked triangle comparison.}
In the free geometric pro-three group on the marked
peripheral generators $x,y$, form
\[
 \Delta=\langle x,y\mid x^3=y^3=(xy)^3=1\rangle_{\rm pro\text{-}3}.
 \tag{E.134.1}
\]
Put $a=yx^{-1}$, $b=xax^{-1}$ and
$c=x^2ax^{-2}$. The kernel of the homomorphism
$\Delta\to\mathbb Z/3$ sending both $x$ and $y$ to
one is generated by $a,b,c$. Rewriting the two
remaining relators gives $abc=1$ and $bac=1$.
Hence $[a,b]=1$, $c=(ab)^{-1}$, and conjugation by
$x$ sends $a$ to $b$ and $b$ to $(ab)^{-1}$.
Conversely these relations define a semidirect
product and imply all three original relators.
These mutually inverse pro-three presentations prove
\[
 \Delta\simeq\mathcal O\rtimes\mathbb Z/3,
 \qquad \mathcal O=\mathbb Z_3[\zeta_3],
 \tag{E.134.2}
\]
where the cyclic generator acts by multiplication by
$\zeta_3$. One may also obtain the same pro-three
presentation by completing
$\mathbb Z[\zeta_3]\rtimes\mathbb Z/3$; its quotients
modulo $3^n$ separate every element.

Since the translation subgroup is abelian,
\[
 \Gamma_i\Delta=(\zeta_3-1)^{i-1}\mathcal O
       \quad(i\ge2),\qquad
 \Gamma_3\Delta=3\mathcal O,\qquad
 \Gamma_5\Delta=9\mathcal O,
 \tag{E.134.3}
\]
using $(\zeta_3-1)^2=-3\zeta_3$.
Thus $\Delta/\Gamma_3\Delta$ is the marked
nonabelian exponent-three group $H_3$, and
\[
 Q_\triangle=\Delta/\Gamma_5\Delta,
 \qquad |Q_\triangle|=243,
 \qquad
 \ker(Q_\triangle\to H_3)=3\mathcal O/9\mathcal O.
 \tag{E.134.4}
\]
The group has class exactly four. It is an actual
quotient of $N_4$, not a comparison with an unrelated
finite group. The normal subgroup defining
$\Delta$ is preserved by Galois, because the three
peripheral conjugacy classes are preserved up to
cyclotomic powers.

To identify the kernel arithmetically, use the actual
pointed cover $Y\to U$ of E.110. Each of its boundary
inertia groups is generated by a conjugate of
$x^3,y^3$, or $(xy)^3$. Passing from $Y$ to its
compactification kills precisely the closed normal
subgroup generated by these boundary inertias.
Therefore the kernel of $\Delta\to H_3$ is the
geometric pro-three fundamental group of $E'$,
canonically $T_3E'$ after choosing its rational
origin. Under (E.134.3) it is $3\mathcal O$, so
the kernel in (E.134.4) is exactly $T_3E'/3T_3E'
=E'[3]$. The point-three section of $U$ lifts to
the rational origin on this pointed cover.
This proves the identification with the actual
Galois action, including its splitting.

\emph{A finite compact class is already a rational Kummer class.}
Let $s_4$ be any global integral class-four lift of
the distinguished class-two section of an E.123 pair.
The latter has the same marked $H_3$ reduction as the
point-three section. Its image in (E.134.4) therefore
defines, after this alignment, a class
\[
                         z\in H^1(\mathbb Q,E'[3]).
 \tag{E.134.5}
\]
Suppose only that $s_4$ matches the local ordinary
point section at two. We claim that $z$ belongs
to the image of $E'(\mathbb Q)/3$.

Every class-three lift of the distinguished section
is point matched at each support prime different
from three, by E.124. At a fresh support prime,
align this class-three point section with the
restriction of $s_4$. Their class-four difference
is an integral weight-four cocycle, whose inertia
is zero by the tame relation. The local point
inertia is a conjugate of $x^{n_j}$ or $y^{n_j}$,
and vanishes in $\Delta$ because $3\mid n_j$.
The compact class (E.134.5) is consequently
unramified there. At every other $q>3$ the
abelian inertia is zero; induction with the
torsion-free relation $(q^i-q)u=0$ shows that
the original integral section is unramified
in every nilpotent class. The reference point-three
section is unramified there as well. Thus $z$ is
unramified at every $q>3$.

The elliptic curve $E'$ has good reduction at all
these places. The unramified subgroup of
$H^1(\mathbb Q_q,E'[3])$ is its local Kummer
subgroup: reduce to the finite residue field and
use $H^1(\mathbb F_q,E')=0$ in the multiplication
sequence. For this vanishing, Frobenius minus the
identity is a nonzero isogeny of $E'$ and is
surjective on its algebraically closed points;
therefore every procyclic one-cocycle is a
coboundary. Hence the image
\[
                         \xi\in H^1(\mathbb Q,E')[3]
\]
of $z$ is locally zero for every $q>3$.
At two the assumed point match gives a rational
point on the pointed cover after its $H_3$
projection, so again $z_2$ is a proper local
Kummer class and $\xi_2=0$. The real group has
no three-primary part.

Global reciprocity for the elliptic Tate pairing
now says that $\xi_3$ annihilates all rational
points; equivalently apply the applicable
localization sequence of
\citet[I.6.26(b)]{CohMilneADT}.
E.112 proves that the rational three-adic
completion maps isomorphically onto
$E'(\mathbb Q_3)^{\wedge3}$. Nondegeneracy of
local duality therefore forces $\xi_3=0$.
E.111 proves $\operatorname{Sha}(E'/\mathbb Q)[3]=0$, so
$\xi=0$. The finite Kummer exact sequence
proves the claim about $z$.
This argument uses one finite coefficient group;
it does not infer a Tate-module class from an
isolated finite lift.

\emph{The support-prime contradiction.}
For a prime $p\equiv-1\pmod3$, cubing is a
bijection on $\mathbb F_p$. Counting by the
$y$-coordinate in $y^2=x^3+9$ gives
$\#E'(\mathbb F_p)=p+1$.
Its rational three-torsion has order exactly
three: it contains $T$, whereas full rational
three-torsion would imply $\mu_3\subset\mathbb F_p$
by the Weil pairing. Thus the three-primary
group is cyclic, of order $3^{v_3(p+1)}$.

At every support prime, the class-two match
already supplies a rational point of the
actual $H_3$ cover over the ordinary local point.
Its compact reduction $R$ obeys the boundary
identity of E.113,
\[
          3R+P_0=-T\quad(A\text{-support}),
          \qquad 3R+P_0=T\quad(B\text{-support}).
 \tag{E.134.6}
\]
For an $A$-support in E.123 the three-primary
group has order three. The first identity
already says $P_0+T=0$ in that group, proving
its compact test.

At a $B$-support $r$, the group is cyclic of
order nine. The second identity implies
$P_0-T\in3E'(\mathbb F_r)^{\wedge3}$.
If the compact test fails, this is a nonzero
element of order three. Both $P_0$ and $T$
then belong to the order-three subgroup, and
their classes in $E'(\mathbb F_r)/3$ are zero.
Since they generate $E'(\mathbb Q)/3$ by E.111,
every rational Kummer class has zero localization
in $H^1(\mathbb Q_r,E'[3])$.

On the other hand, any compact boundary lift
$R$ in (E.134.6) has $3R=-(P_0-T)\ne0$
in the three-primary group. Its three-primary
component has order nine, so its class modulo
three is nonzero. If $s_4$ also matches the
ordinary point at $r$, its finite compact
class $z_r$ is exactly this nonzero Kummer
class. This contradicts the rational Kummer
conclusion proved above from the match at two.

The construction in E.129 imposes precisely
such a nonzero $P_0-T$ at its first $B$-support.
It is therefore excluded at this fourth
matching layer. The finite quotient and the
two witnessing support places are explicit,
and the coefficient place three was used only
for duality and the proved localization
isomorphism, not for a point-section assumption.
\end{proof}

\begin{result}{E.135}{Proved}{The finite-field point--cusp question is an integral Frobenius-path test}
Let $k=\mathbb F_q$, with $q$ prime to three, let
$t\in\mathbb P^1(k)\setminus\{0,1,\infty\}$, and fix a
rational tangential section at one of the three cusps.
Write $\Pi_3$ for the free geometric pro-three group,
$N_r=\Pi_3/\Gamma_{r+1}\Pi_3$, and $U_r$ for its
rational unipotent completion over $\mathbb Q_3$.
Use arithmetic Frobenius $F$, acting relative to the
tangential section, and let $z_t$ be the point-section
cocycle evaluated at $F$.

There is a unique compatible sequence
$g_r\in U_r(\mathbb Q_3)$ satisfying
\[
                 z_t=g_r^{-1}F(g_r)\quad\text{in }U_r(\mathbb Q_3).
 \tag{E.135.1}
\]
The point and tangential section classes in the full
geometrically pro-three arithmetic fundamental group
coincide if and only if $g_r\in N_r$ for every $r$.

If $q\equiv-1\pmod3$ and all earlier coordinates are
integral, a first failure of integrality can only occur
at an even weight $2k$. Its exact correction denominator
has valuation
\[
                     v_3(q^{2k}-1)=v_3(q^2-1)+v_3(k).
 \tag{E.135.2}
\]
The elliptic detector E.192 now proves that this
integrality test always fails.  If $N$ is the least
integer with $3^N>2q+1$, the first failure has weight
at most $3N-2$.  For $q\equiv-1\pmod3$ it is even,
and is therefore at most the largest even integer
not exceeding $3N-2$.
This local separation does not resolve the main
global ghost problem.
\end{result}
\begin{proof}
The integral Hall coordinates identify $N_r$ with a
compact torsion-free nilpotent three-adic group and
embed it in $U_r(\mathbb Q_3)$. On its weight-$i$
central quotient Frobenius acts by $q^i$.
At weight one, (E.135.1) is the linear equation
$(F-1)v=\log(z_t)$, whose coefficient $q-1$ is
nonzero in $\mathbb Q_3$.
Suppose the unique solution has been obtained in
class $r-1$, and lift it arbitrarily to $U_r$.
Its discrepancy from (E.135.1) belongs to the
central weight-$r$ vector group. Multiplying the
chosen lift by a central vector $v$ changes
$g^{-1}F(g)$ by $(q^r-1)v$.
This coefficient is again nonzero in $\mathbb Q_3$,
so there is exactly one correction. Induction proves
existence, uniqueness, and compatibility of the
rational solutions.

An integral geometric conjugator necessarily maps
to these unique $g_r$. Conversely, if each $g_r$
is integral, the corresponding fibres of
$\Pi_3\to N_r$ are nonempty nested compact sets.
Their intersection supplies an element of $\Pi_3$
satisfying the equation. Equality on arithmetic
Frobenius implies equality on all of $G_k$ by
continuity, because its integer powers are dense.
Here $\Pi_3\simeq\varprojlim_rN_r$: every finite
three-group quotient is nilpotent, so the lower
central series has trivial intersection, while
compactness gives surjectivity to the inverse limit.

If the solution at class $r-1$ is integral, it has
an integral lift to $N_r$. The remaining discrepancy
is therefore an integral central vector. Changing
the lift changes that vector by a multiple of
$q^r-1$. Integrality at the next step is precisely
its vanishing in
\[
                   \bigl(\mathbb Z_3/(q^r-1)\bigr)^{b_r}.
\]
For $q\equiv-1\pmod3$, this group is zero for
odd $r$. For $r=2k$, the elementary lifting-the-exponent
identity gives (E.135.2). A nonzero discrepancy is
detected in a finite three-group quotient, by compact
separation of the integral conjugacy class.
The reduction itself supplies no bound on the first
detecting weight.  The separate representation proof
E.192 does: its geometric image modulo $3^N$ has order
at most $3^{3N-2}$, hence nilpotency class at most
$3N-2$.  An integral conjugator in that class would
map to a conjugator in the detecting finite image,
contradicting its two distinct Frobenius traces.
This proves the stated bound; the already proved
odd-weight vanishing gives its even refinement.

Two relevant primary sources were checked for a
stronger separation conclusion.
\citet[Proposition~1.8, pp.~144--147, and Corollary~1.11, pp.~147--148]{AuditSaidiTamagawa2009}
separate full cusp decomposition groups from other
decomposition-group collisions. A tangential section
is a procyclic complement to cusp inertia, so equality
of that complement with an interior point section
does not identify the full decomposition groups.
The source's finiteness assertion for its exceptional
set also requires a cofinite set of coefficient primes.
\citet[Theorems~A--C, author pp.~2--4]{AuditWakabayashi2012}
reconstruct cuspidalizations from specified
Frobenius-preserving isomorphisms. Those hypotheses
do not themselves supply the integrality failure
needed in (E.135.1). We consequently use neither
source as a point--cusp separation theorem.
The later proof E.192 instead uses a smooth elliptic
family with rational three-torsion and an actual
split-multiplicative cusp representation.
\end{proof}

\begin{result}{E.136}{Proved}{Unbounded degree and genus do not separate the finite-field point and cusp}
Let $q\equiv-1\pmod3$, and put $m_n=3^n$.
The Fermat covers of the tripod with compactifications
\[
 C_n:\quad X^{m_n}+Y^{m_n}=Z^{m_n},
 \qquad C_n\longrightarrow\mathbb P^1,
 \qquad t=(X/Z)^{m_n}
\]
have geometric degree $3^{2n}$ and genus
$(3^n-1)(3^n-2)/2$. Nevertheless
\[
                         \#C_n(\mathbb F_q)=q+1
                         \quad\text{for every }n.
\]
Every interior $t\in U(\mathbb F_q)$ has exactly one
compatible rational lift throughout this tower, and
each cusp also has one compatible rational lift in
the compactifications.

Thus increasing the degree and genus of a tower of
three-covers while retaining rational lifts of a point
and cusp does not itself contradict the Hasse--Weil
bounds. This particular tower is the geometric
abelian pro-three tower; no nonabelian point--cusp
equality is asserted.
\end{result}
\begin{proof}
Since $\gcd(3,q-1)=1$, every map
$u\mapsto u^{3^n}$ is a permutation of
$\mathbb F_q$, including zero. For each affine
$X/Z\in\mathbb F_q$ there is exactly one
$Y/Z\in\mathbb F_q$ on $C_n$, giving $q$
affine points. At infinity the equation
$(X/Y)^{3^n}=-1$ has exactly one solution.
This proves the point count.

For a fixed interior $t$, both equations
$(X/Z)^{3^n}=t$ and $(Y/Z)^{3^n}=1-t$
have unique solutions in $\mathbb F_q$.
The maps $C_{n+1}\to C_n$ cube every coordinate,
so uniqueness makes these lifts compatible.
The points $(0:1:1),(1:0:1),(1:-1:0)$ give
the three compatible cusp lifts.

Over an algebraic closure the deck group on the
open cover is $(\mu_{3^n})^2$, of order $3^{2n}$.
The plane curve is smooth because its degree is
prime to the characteristic; the plane genus
formula gives the stated genus. These are exactly
the Kummer covers defining the abelian geometric
quotients modulo $3^n$. Equivalently, their
section classes cannot distinguish the points
from a tangent because
$H^1(G_{\mathbb F_q},\mathbb Z_3(1))
=\mathbb Z_3/(q-1)=0$.
Finally $q+1$ lies at the centre of every
Hasse--Weil interval
$[q+1-2g_n\sqrt q,q+1+2g_n\sqrt q]$.
The counts are therefore compatible with all
those bounds even as $g_n$ tends to infinity.
\end{proof}


\begin{result}{E.137}{Proved}{The all-prime class-two support gate has a two-place failure witness}
Let $(A,B)$ be an ordinary completed pair and let
\[
 \Sigma=\{q:(v_q(A),v_q(B))\ne(0,0)\}.
\]
Fix one global integral class-two section class $s$ above the pair.
For each $q\in\Sigma$, let $D_q$ be the difference between its
local class and the prescribed ordinary point class, with the sign
chosen so that $s+k\eta$ matches precisely when $k\eta|_{G_q}=D_q$.
Only coefficient components $\ell\ne q$ are included in this
comparison. After compatible Tate coordinates are fixed,
\[
 D_q\in\prod_{\ell\ne q}H^1(G_{\mathbb Q_q},\mathbb Z_\ell(2))
       \simeq \mathbb Z/(q^2-1).
 \tag{E.137.1}
\]
In particular, only the finitely many primes dividing $q^2-1$
can contribute a local difference.

For $q\ge5$, localization of the global twist group is the injection
\[
 \mathbb Z/24\longrightarrow\mathbb Z/(q^2-1),\qquad
 k\longmapsto k\frac{q^2-1}{24}.
 \tag{E.137.2}
\]
Thus matching at a single support prime $q_0\ge5$ determines
at most one of the twenty-four global classes. If $\Sigma$
contains such a prime, the set of global class-two classes matched
at every support prime in every off-diagonal coefficient component
has at most one member. If this set is empty, at most two support
primes already witness its emptiness. No effective bound on the
sizes of those primes is asserted.

Equivalently, each $D_q$ for $q\ge5$ must lie in the order-$24$
subgroup in (E.137.2), its unique preimages must all agree, and
any support conditions at $q=2,3$ must agree with their reductions
modulo three and eight, respectively. If these conditions hold,
B.25 gives a global integral class-three lift of the unique matched
class-two section. This does not assert that the class-three lift
remains matched, or that any higher lift exists.

Every ordinary full pronilpotent survivor satisfies this gate by
A.31. The gate concerns support primes with $q\ne\ell$ only.
Matching at the coefficient prime and at zero-valuation places
has not been deduced from full pronilpotent survival.
The existence of an infinite-support ordinary pair passing this
all-prime gate remains unresolved in this work.
\end{result}
\begin{proof}
By A.25--A.27 the global class-two fibre is a nonempty torsor
under $H^1(G_{\mathbb Q},\widehat{\mathbb Z}(2))
=\mathbb Z/24$, with generator $\eta$.
At $q\ne\ell$, inertia acts trivially on $T=\mathbb Z_\ell(2)$.
Its weight-one-shifted tame cohomology has no Frobenius invariants:
\[
 H^1(I_q,T)^{G_{\mathbb F_q}}
 =\mathbb Z_\ell(1)^{G_{\mathbb F_q}}=0,
\]
since multiplication by $q-1$ is injective on $\mathbb Z_\ell$.
Wild inertia contributes nothing. Inflation and the procyclic
calculation consequently give
\[
 H^1(G_{\mathbb Q_q},T)
 =H^1(G_{\mathbb F_q},T)
 =\mathbb Z_\ell/(q^2-1).
 \tag{E.137.3}
\]
The local fibre above the fixed abelian class is a torsor under
this group: align the literal abelian cocycles, then take the
central difference. There is no additional stabilizer quotient,
because $H^0(G_{\mathbb Q_q},\mathbb Z_\ell(1)^2)=0$.
Any two choices of the central alignment differ by a coboundary.
The local point section provides a distinguished element of this
fibre. Taking the product over $\ell\ne q$ proves (E.137.1),
since $q$ does not divide $q^2-1$.

The cocycle $\eta=(\chi^2-1)/24$ is unramified at $q$ in
every off-diagonal coefficient component, and its value on
arithmetic Frobenius is $(q^2-1)/24$. For a prime $q\ge5$,
the integer $q^2-1$ is divisible by $24$. Its indicated class
has exact order $24$, proving (E.137.2). At $q=2$ only the
three-primary local group is nonzero; its generator value is
$1/8$ modulo three, so localization factors through an isomorphism
from the global order-three factor. At $q=3$ the corresponding
value is $1/3$ modulo eight, giving an isomorphism from the global
order-eight factor. These statements include the units and do
not use an invalid integer division at the small primes.

Choose a support prime $q_0\ge5$. If no global class matches
there, that single place proves emptiness. Otherwise there is
exactly one such global class, by injectivity. If it does not
match all the support places, choose a support prime $q_1$
where it fails. Every other global class already fails at $q_0$,
so $\{q_0,q_1\}$ is a failure witness. Each of these two
place conditions uses only finitely many coefficient primes
by (E.137.3). An infinite support set necessarily contains a
prime at least five, so the assertion applies to every ghost
under consideration.

If the matching conditions all hold, apply B.25 with $r=2$
separately at every coefficient prime. Its strict residue argument
includes all three ordinary valuation patterns and gives a global
integral class-three lift in each component, because the fixed
ramification weight-three second cohomology is zero. Their product
is a global integral class-three lift. No simultaneous local
matching condition at the new layer follows from that existence
argument. Finally A.31 is exactly the stated off-diagonal support
necessity; it supplies no condition at the omitted local places.
\end{proof}


\begin{result}{E.138}{Proved}{Matching at the support prime five imposes no abelian rigidity}
Every ordinary completed pair whose support contains five has exactly
one global integral class-two section class which matches its prescribed
ordinary point at five in every coefficient component different from
five. This condition imposes no additional restriction on the abelian
pair. In particular there are infinite-support ordinary examples
satisfying it, furnished by E.125.
\end{result}
\begin{proof}
At five the entire off-diagonal local difference group of E.137 has
order $5^2-1=24$. Its localization map from the global twist group is
\[
 \mathbb Z/24\longrightarrow\mathbb Z/24,\qquad
 k\longmapsto k(5^2-1)/24=k.
\]
It is an isomorphism. Every local difference is therefore corrected
by exactly one global twist, independently of the abelian pair.
The E.125 examples have valuation pair $(1,0)$ at five and infinite
support, so they give the asserted counterexamples to any inference
of finite support from this single local matching condition.
Neither matching in coefficient five at this place nor matching at
another place follows from the argument.
\end{proof}


\begin{result}{E.139}{Proved}{The relative Bloch component, its harmonic cancellation, and a nonzero ghost example}
Let $q\ge5$ be a support prime of an ordinary pair, and let
$\ell$ be an odd prime dividing $q+1$. The $\ell$-component
of the local difference in E.137 has a natural interpretation
\[
 D_{q,\ell}\in H^1(G_{\mathbb F_q},\mathbb Z_\ell(2))
       \simeq K_3(\mathbb F_q)\{\ell\}
       \simeq \mathcal B(\mathbb F_q)\{\ell\},
 \tag{E.139.1}
\]
where $\mathcal B$ is Suslin's Bloch group. For $\ell\ge5$
there is no global twist ambiguity. For $\ell=3$, the ambiguity
is exactly the order-three subgroup generated by the localization
of $\eta$. The resulting quotient class is therefore an invariant
of the ordinary completed pair. This is a relative class defined
by its global central primitive and the local point comparison;
it has not been identified with the single symbol of the point's
reduction at a cusp.

There are infinite-support ordinary pairs for which this relative
Bloch quotient is nonzero. More precisely, a modification of the
E.123 recursion has a $B$-support prime $r$ with
$v_3(r^2-1)=2$ at which $D_{r,3}$ has order nine for every
global integral class-two section. Thus none of the twenty-four
global classes matches there, although the prescribed ordinary
point reduces to the cusp one.

There is also an exact harmonic cancellation for the rational
reference point $t=1/2$: for every $\ell\ge5$ and every
$q\ne2,\ell$ with $\ell\nmid q-1$, its integral class-two
local section equals the unit tangential section at zero.
In particular its $q+1$ Bloch component vanishes. This reference
point calculation does not assert the same cancellation after
its two global Kummer coordinates have been changed to a ghost.
\end{result}
\begin{proof}
The inflation isomorphism in E.137 identifies the local difference
with finite-field cohomology. The finite-field calculations give
$K_2(\mathbb F_q)=0$ and
$K_3(\mathbb F_q)=\mathbb Z/(q^2-1)$; the coefficient Chern
comparison is an isomorphism with weight-two first cohomology.
For the exact comparison map, see
\citet[Theorem~2.10, author p.~4]{AuditHutchinson2024};
the compatible local diagram is written in
\citet[Lemma~4.1, pp.~403--404]{AuditCGZ2023}.
Taking inverse coefficient limits gives the first isomorphism
in (E.139.1). Suslin's exact sequence has kernel the extension
of the root-of-unity torsion of order $q-1$ by a group of order
two, and its Bloch quotient has order $(q+1)/2$
\citep[VI.5.1.1--5.2]{CohWeibelKbook}.
Since an odd $\ell$ dividing $q+1$ does not divide $q-1$,
it induces the second isomorphism. No such isomorphism with the
Bloch group is claimed for the $q-1$ primary components or at two.

The global ambiguity is exactly A.26. It is zero for
$\ell\ge5$ and has order three at $\ell=3$; E.137 proves
that its local image has that same order. This also explains
why adding globally integral cyclotomic or Bloch cocycles while
holding $(A,B)$ fixed cannot repair the relative quotient class:
every such addition is already one of these global twists.
The finite-coefficient Bloch group is larger. Its extra quotient
is $K_2(F)[\ell^n]$, with the coefficient Bockstein comparison
made explicit in \citet[Theorem~1.7, p.~387, and \S6,
pp.~414--416]{AuditCGZ2023}. A finite correction must consequently
not be promoted to a compatible integral correction without
checking this extra condition.

We now construct the nonzero example, using only the previously
proved finite-field compatibility in E.123 and E.128. Run E.123,
with $n_j=(j+7)!$, and retain all its congruences, local roots,
and marked Frobenius prescriptions except at one chosen new
$B$-support prime $r$. Retain the shell
\[
 r\equiv-1\pmod9,\qquad r\not\equiv-1\pmod{27},
\]
and retain its abelian marked coordinates
$\alpha=\beta=0\pmod9$. Replace its central point-three
$H_9$ Frobenius prescription by
\[
                 \gamma=\eta(F_r)+1\pmod9.
 \tag{E.139.2}
\]
As always, $\chi(F_r)$ is retained modulo $27$ in order to
specify $\eta(F_r)$ modulo nine.
The point-three $H_9$ image contains its full center. The cyclic
intersection with the current radical extension in E.123 is
its abelian ninth-root line. Multiplying the prescribed element
by the central generator changes neither that line nor any
cyclotomic or root condition. The fixed-root compatibility of
E.128 therefore supplies the required prime. All subsequent
choices are made as in E.123. Its convergence and ordinarity
proofs depend on the unchanged root and congruence conditions,
so the final pair is still ordinary and has infinite support.

The integral class-two normalization in E.123 also remains valid:
its proof uses divisibility of every new inertia exponent by nine
and the unchanged old-place comparisons, not the value chosen
for the new Frobenius center. Thus choose the global class with
$h\equiv h_{\rm point\,3}\pmod9$.
At $r$, put $v=\alpha/(r-1)$. Since $r-1$ is a three-adic
unit, $v=0\pmod9$, and the exact boundary mismatch is
\[
       \gamma+\beta v-\eta(F_r)=1\pmod9.
 \tag{E.139.3}
\]
The denominator $r^2-1$ has valuation two. Accordingly its class
in $H^1(G_r,\mathbb Z_3(2))$ has exact order nine.
A global twist adds a multiple of $\eta(F_r)$, whose valuation
is one, so it cannot change (E.139.3) into zero even modulo
three. The other coefficient twists have no effect on this
component. This excludes all twenty-four global classes at
that single support place. Since $r\equiv-1\pmod3$, (E.139.1)
identifies this order-nine class with the three-primary Bloch
component; its quotient by the global order-three subgroup is
nonzero. No claim of a class-three lifting obstruction follows
from this calculation: multiplication by the large ordinary
inertia exponent can kill that later local obstruction.

Finally consider the harmonic reference point. In the integral
Heisenberg group at an odd coefficient prime, replace the Hall
central coordinate $h$ by $z=h-ab/2$. The group law is
\[
 (a,b,z)(a',b',z')=
 (a+a',b+b',z+z'+(ab'-ba')/2),
\]
and the actual action is $(a,b,z)\mapsto(\chi a,\chi b,\chi^2z)$.
If $a=\kappa_\ell(1/2)$ is a literal Kummer cocycle, then
$g\mapsto(a(g),a(g),0)$ is a global class-two cocycle: the
alternating cross term is identically zero. Its abelianization
is that of the actual point $1/2$. For $\ell\ge5$, A.26 makes
the global class-two fibre a singleton, so this is the actual
point class. At the indicated prime $q$, the Kummer class is
unramified and $q-1$ is an $\ell$-adic unit. Write
$a(g)=(\chi(g)-1)u$ locally. The diagonal element $(u,u,0)$
then has coboundary $(a,a,0)$, proving the asserted integral
local equality with the tangent. A ghost has two independently
changed Kummer coordinates; the vanished alternating term in
this calculation is consequently no longer available. Neither
the Chern comparison nor the harmonic calculation removes that
missing cross-term control.

The exact finite script \path{code/ghost_audit_bloch.py}
checks the cyclic-dilogarithm five-term relation in five residue
fields, the harmonic values, and both central-shift shells modulo
$27$; its output is saved in \path{code/ghost_audit_bloch.out}.
These are finite checks of the formulas and their scope, not a
substitute for the preceding Chebotarev existence proof.
\end{proof}
